\documentclass[12pt]{article}

% === Packages ===
\usepackage[utf8]{inputenc}
\usepackage[T1]{fontenc}
\usepackage{lmodern}
\usepackage[margin=1in]{geometry}
\usepackage{amsmath,amssymb}
\usepackage{graphicx}
\usepackage{booktabs}
\usepackage{tabularx}
\usepackage{xcolor}
\usepackage[round]{natbib}
\usepackage[colorlinks=true,linkcolor=blue!60!black,citecolor=blue!60!black,urlcolor=blue!60!black]{hyperref}
\usepackage{microtype}
\usepackage{setspace}
\usepackage{parskip}
\usepackage{placeins}
\usepackage{float}
\usepackage{enumitem}

% === Formatting ===
\onehalfspacing

% === Title ===
\title{The Four-Model Theory of Consciousness:\\A Simulation-Based Framework Unifying the Hard Problem, Binding, and Altered States}
\author{Matthias Gruber\\
\textit{Independent researcher}\\
ORCID: \href{https://orcid.org/0009-0005-9697-1665}{0009-0005-9697-1665}\\
\texttt{matthias@matthiasgruber.com}}
\date{}

% === Document ===
\begin{document}
\maketitle

\begin{abstract}\singlespacing
No existing theory of consciousness simultaneously satisfies the eight core requirements a complete theory must meet: the Hard Problem, the Explanatory Gap, the Boundary Problem, the Structure of Experience, Unity and Binding, Combination and Emergence, the Causal Role, and the Meta-Problem. This paper presents the Four-Model Theory, in which consciousness is constituted by ongoing self-simulation across four nested models arranged along two axes --- scope (world vs. self) and mode (implicit vs. explicit). The implicit models (Implicit World Model, Implicit Self Model) are substrate-level, learned, and non-conscious. The explicit models (Explicit World Model, Explicit Self Model) are virtual, transient, and phenomenal --- they are the simulation in which experience occurs. The theory's central claim is that qualia are constitutive properties of the computational level --- digital constructs that exist at the level of the running computation but are incoherent at the substrate level, just as a spreadsheet cell's value is incoherent at the transistor level. The theory addresses the Hard Problem by revealing a category error --- a level confusion that seeks phenomenal properties at the substrate level where they categorically do not exist. Self-referential closure explains why this specific computational process has experience when a weather simulation does not: the system's model includes a model of itself, collapsing the inside/outside distinction and making experience constitutive rather than additional. Combined with a criticality requirement (the substrate must operate at the edge of chaos), the theory derives diverse phenomena from five principles: criticality, virtual qualia, a redirectable Explicit Self Model, variable implicit-explicit permeability, and simulation forking. These principles unify psychedelic phenomenology, anesthetic mechanisms, dream states, split-brain phenomena, dissociative identity disorder, and animal consciousness. A systematic comparison shows the theory addresses all eight requirements. Unusually for a consciousness theory, the framework has substantial empirical grounding: five claims that follow from its core axioms --- developed over the preceding decade and published in 2015 --- have since been corroborated by independent research groups with no connection to the theory, including the anesthetic-criticality convergence \citep{Casali2013,HengenShew2025,AlgomShriki2026}, sleep-dependent criticality restoration \citep{Bhatt2024}, sleep onset as bifurcation \citep{Li2025}, and split-brain holographic degradation \citep{Pinto2017}. Four novel predictions remain untested --- including that psychedelics should alleviate anosognosia and that ego dissolution content is controllable via sensory input --- predictions no competing theory generates.
\end{abstract}

\noindent\textbf{Keywords}: consciousness, hard problem, self-model, simulation, qualia, criticality, binding problem, altered states, psychedelics, substrate independence




%=============================================================================
\section{Introduction}
%=============================================================================



\subsection{The State of Consciousness Science}


After three decades of intensive scientific investigation, consciousness research has not converged on a dominant paradigm \citep{Kuhn1962}. The field possesses no agreed-upon methodology for linking subjective experience to objective measurement, and no theory that commands broad assent. Multiple frameworks compete for explanatory primacy --- Integrated Information Theory (IIT; \citealp{Tononi2004,Albantakis2023}), Global Neuronal Workspace (GNW; \citealp{Baars1988,Dehaene2011}), Higher-Order Theories (HOT; \citealp{Rosenthal2005,Lau2011}), Predictive Processing (PP; \citealp{Friston2010,Seth2021}), Attention Schema Theory (AST; \citealp{Graziano2013}), Recurrent Processing Theory (RPT; \citealp{Lamme2006,Lamme2010}), and others --- yet none has established decisive empirical or theoretical superiority over its rivals.

Recent developments illustrate the difficulty. The COGITATE adversarial collaboration published equivocal results in \emph{Nature} \citep[protocol:][]{COGITATE2025,Melloni2023}: neither IIT nor GNW was fully confirmed, and the data showed stronger posterior cortical involvement than GNW predicted while leaving IIT's exclusion postulate unresolved. A letter signed by over 100 researchers declared IIT pseudoscientific \citep{IITConcerned2025}, provoking fierce rebuttals \citep{TononiAlbantakis2025} and methodological commentary \citep{GomezMarin2025}. Reviews published in 2024-2025 have asked whether the field is making genuine progress or merely accumulating incompatible frameworks \citep{SethBayne2022,KirkebyHinrup2025b,WagnerAltendorf2024}.

Moreover, the field's leading theories have limited track records of novel empirical prediction: most empirical work in consciousness science involves post-hoc correlation of neural measures with existing theories, not prospective prediction from theoretical first principles subsequently confirmed by independent groups. This paper argues that the impasse persists because no existing theory simultaneously addresses all of the fundamental requirements that a complete theory of consciousness must meet. Each theory excels on some requirements but remains silent on, or weak against, others. The field needs a framework that addresses the full set of challenges --- not just neural correlates or access conditions, but the deep philosophical problems that make consciousness uniquely difficult.


\subsection{What Would Count as Progress?}


This paper proposes that any theory claiming to provide a comprehensive account of consciousness must address eight core requirements, drawn from the philosophical and scientific literature. These requirements are not novel; each has been identified by previous authors as a central challenge. What is novel is the demand that a single theory address all eight simultaneously:

\begin{enumerate}
\item \textbf{The Hard Problem} \citep{Chalmers1995} --- Why does physical processing give rise to subjective experience?
\item \textbf{The Explanatory Gap} \citep{Levine1983} --- Why does the explanation of neural correlates feel incomplete?
\item \textbf{The Boundary Problem} \citep{Bayne2010,Tononi2004} --- Where does the conscious system end?
\item \textbf{The Structure of Experience} \citep{Nagel1974} --- How does physical processing produce richly structured experience?
\item \textbf{Unity and Binding} \citep{Treisman1980,Revonsuo1999} --- How are distributed processes unified into coherent experience?
\item \textbf{Combination and Emergence} \citep{Chalmers2016} --- How do non-conscious elements combine to produce consciousness?
\item \textbf{The Causal Role} \citep{Jackson1982} --- Does consciousness do anything?
\item \textbf{The Meta-Problem} \citep{Chalmers2018} --- Why do we think there is a hard problem?
\end{enumerate}

Section 2 develops each requirement in detail and surveys how existing theories fare against them. The remainder of the paper presents a theory --- the Four-Model Theory --- that addresses all eight.


\subsection{Overview and Historical Context}


The Four-Model Theory was originally published in German as \emph{Die Emergenz des Bewusstseins} \citep{Gruber2015} and has been refined through a structured adversarial challenge process in 2026. The theory takes Dennett's Multiple Drafts Model \citep{Dennett1991} and Metzinger's Self-Model Theory of Subjectivity \citep{Metzinger2003,Metzinger2009} as its starting point --- but as the 2015 monograph noted, the combination of MDM and SMT leaves several things unexplained. The theory's substance lies in what it builds beyond that starting point. It is substrate-independent: the six-layer mammalian cortex is understood as an evolutionary implementation, not a requirement. Specifically, the theory adds five elements that neither MDM, SMT, nor any other existing framework provides individually or in combination:

\begin{enumerate}
\item \textbf{The cortical automaton concept.} Consciousness is identified with a specific class of physical dynamics: the high-dimensional activity pattern that constitutes a cellular automaton operating at criticality on the cortical substrate. This differs from IIT, which identifies consciousness with a mathematical measure (integrated information, $\Phi$) that in principle applies to any substrate meeting the measure's threshold --- FMT's identification is with a specific \emph{kind} of dynamics (CA at criticality), not with a substrate-neutral quantity. The identification is not metaphorical; it carries measurable signatures \citep{Gruber2015}. No other theory specifies this particular dynamical interpretation.
\item \textbf{A two-level ontology that addresses the Hard Problem.} Qualia are constitutive properties of the computational level --- real at that level but categorically incoherent at the substrate level. This is distinct from Dennett's illusionism (which denies qualia are real) and from Metzinger's transparency thesis (which explains why the self-model is not experienced as a model but does not claim to address the Hard Problem). It constitutes a third position: qualia are real, level-specific, and the Hard Problem rests on a category error.
\item \textbf{A criticality requirement derived from computational theory.} The edge-of-chaos requirement was derived from Wolfram's classification of cellular automata \citep{Wolfram2002} --- in 2015, a decade before the empirical confirmation by the ConCrit framework \citep{AlgomShriki2026} and \citet{HengenShew2025} meta-analysis of over 140 datasets. No other consciousness theory derived this requirement from first principles prior to its empirical confirmation.
\item \textbf{Variable implicit-explicit permeability as a unified mechanism for altered states.} A single principle --- modulation of the boundary between substrate-level (implicit) and simulation-level (explicit) processing --- generates the phenomenology of psychedelics, dreams, meditation, anesthesia, and anosognosia. Neither MDM nor SMT provides a mechanism that unifies these phenomena under one parameter. The closest competitor is the REBUS model \citep{CarhartHarris2019}, which unifies altered states under relaxed precision weighting of prediction error; however, REBUS operates at the level of hierarchical message passing, whereas FMT's permeability operates at the structural boundary between substrate and simulation --- a level that is architecturally prior and from which the REBUS mechanism is derivable as one implementation.
\item \textbf{A graduated consciousness hierarchy via recursive self-modeling depth.} The theory defines a principled continuum from simple consciousness (a system that models itself) through triply extended consciousness (a system that models its modeling of its modeling of itself). This provides graded rather than binary criteria for attributing consciousness across species and artificial systems, without ad hoc thresholds.
\end{enumerate}

The theory belongs to the self-modeling tradition in consciousness studies, which holds that consciousness is constituted by --- rather than merely correlated with --- the system's model of itself. This tradition includes Metzinger's Self-Model Theory of Subjectivity \citep{Metzinger2003,Metzinger2009}, which analyzed the phenomenal self-model as a transparent representational structure; Damasio's multi-level self framework \citep{Damasio1999,Damasio2010}, which grounded consciousness in progressively elaborated self-representations; Hofstadter's strange-loop account \citep{Hofstadter2007}, which identified the self with self-referential symbolic structures; Graziano's Attention Schema Theory \citep{Graziano2013}, which modeled subjective awareness as the brain's schema of its own attention; and Seth's interoceptive inference framework \citep{Seth2021,SethBayne2022}, which cast the experienced self as a controlled hallucination grounded in body-state prediction. The Four-Model Theory extends this tradition by specifying the minimal architecture required for consciousness-constituting self-simulation (four model kinds along two axes), adding a physical prerequisite (criticality), and developing a two-level ontology (real/virtual) that addresses the Hard Problem.

The theory proposes that consciousness consists of an ongoing self-referential computation running on an implicit (substrate-level) knowledge base. Qualia --- the subjective ``feel'' of experience --- are constitutive properties of the computational level: they exist at the level of the running computation but are incoherent at the substrate level. This two-level ontology addresses the Hard Problem by showing it rests on a category error --- a level confusion.

Combined with a criticality requirement derived independently from Wolfram's computational framework \citep{Wolfram2002}, the theory generates four novel testable predictions and unifies phenomena across psychopharmacology, clinical neurology, sleep science, and comparative cognition under five principles. Five claims that follow from the theory's axioms --- established in 2015 --- have since been corroborated by subsequent empirical research from independent groups, a track record unusual among consciousness theories (Section 8.1).

The theory is offered as one model among several, contributing to humanity's collective search for an adequate account of consciousness. Every model carries inherent modeling error; the present theory is no exception. It is intended to complement existing frameworks --- extending where they are incomplete, not displacing where they succeed.

The paper proceeds as follows. Section 2 establishes the eight requirements. Section 3 presents the Four-Model Theory. Sections 4 and 5 develop the philosophical commitments and the binding/criticality framework. Section 6 demonstrates the theory's explanatory range. Section 7 provides a systematic comparative analysis against major competitors. Section 8 presents the empirical convergence evidence and four predictions. Sections 9-11 address open questions, implications, and conclusions.



%=============================================================================
\section{Eight Requirements for a Theory of Consciousness}
%=============================================================================


This section develops each of the eight requirements and briefly notes how current theories address or fail to address them. A detailed theory-by-theory comparison follows in Section 7; the purpose here is to establish the evaluative framework.


\subsection{The Hard Problem}


Chalmers (1995, 1996) formulated the Hard Problem as the question of why physical processing is accompanied by subjective experience. We can explain all the \emph{functions} of consciousness --- discrimination, integration, reporting, attention --- without explaining why there is ``something it is like'' \citep{Nagel1974} to undergo these processes. The explanatory challenge is not about identifying neural correlates but about understanding why correlates are accompanied by phenomenality at all.

Most neuroscientific theories of consciousness (GNW, RPT, PP) focus on functional or mechanistic accounts and remain explicitly or implicitly silent on the Hard Problem. IIT attempts to address it by defining consciousness as intrinsic causal power ($\Phi$), treating experience as identical to integrated information --- but this requires accepting panpsychist commitments that many find problematic \citep{Aaronson2014,Doerig2019}. HOT and AST offer deflationary accounts that explain \emph{why we report} having phenomenal experience but leave open whether they have truly addressed the phenomenality itself. Illusionism \citep{Dennett1991,Frankish2016} dissolves the problem by denying that qualia as traditionally conceived exist --- a position that remains deeply controversial.


\subsection{The Explanatory Gap}


\citet{Levine1983} identified the Explanatory Gap as a distinct problem from the Hard Problem: even if we identify every neural correlate of every conscious state, the explanation seems to leave something out. The gap is between the third-person description (neural firing patterns) and the first-person reality (what the experience is like). Block (1995, 2007) further refined this as the distinction between access consciousness (the functional role) and phenomenal consciousness (the subjective feel).

The Explanatory Gap is often treated as a restatement of the Hard Problem, but it has a distinct character: it is about the \emph{form} of explanation rather than the \emph{existence} of the phenomenon. A theory that addresses the Hard Problem should simultaneously close the Explanatory Gap.


\subsection{The Boundary Problem}


Where does the conscious system begin and end? Within the brain, only some processing is conscious at any given moment. Between organisms, it is unclear where to draw the line. The Boundary Problem asks for a principled account of what delineates conscious from non-conscious processing \citep[see][]{Bayne2010,Tononi2004}.

IIT provides the strongest existing treatment of this requirement through its exclusion postulate: the system with maximum $\Phi$ defines the boundary of consciousness. However, the computational intractability of calculating $\Phi$ limits its practical application. GNW defines conscious access in terms of global broadcasting, but the boundary between broadcast and non-broadcast content is not always sharp. PP uses Markov blankets \citep{Friston2010} but has been criticized for being too liberal in its boundary-setting \citep{Bruineberg2022}.


\subsection{The Structure of Experience}


Conscious experience is not a homogeneous blob --- it has rich spatial, temporal, modal, and qualitative structure. A visual scene has colors, shapes, depths, and textures; an auditory experience has pitches, timbres, and spatial locations; emotional experience has valence, intensity, and phenomenal character. Any complete theory must explain how physical processing generates this structured phenomenology.

IIT's qualia space provides a mathematical treatment of experiential structure, arguably its greatest strength. PP's generative models are inherently structured, providing a natural account of the richness of perceptual experience. GNW and HOT are weaker here, offering accounts of \emph{when} content becomes conscious but less about \emph{why} it has the particular structure it does.


\subsection{Unity and Binding}


The Binding Problem \citep{Treisman1996,Revonsuo1999} asks how distributed neural processes --- occurring in different brain regions, at different timescales, in different modalities --- are unified into a single coherent conscious experience. I see a red ball: ``red,'' ``round,'' ``ball,'' ``there,'' and ``now'' are processed in different cortical areas, yet I experience a unified percept.

Proposed solutions range from temporal synchrony \citep{Gray1989,Singer1995,Fries2005,Fries2015} to integrated information \citep{Tononi2004} to global broadcasting \citep{Baars1988}. None is universally accepted. The binding problem remains, alongside the Hard Problem, one of the deepest unresolved challenges.


\subsection{Combination and Emergence}


How do non-conscious elements combine to produce consciousness? For panpsychist theories (IIT in its strong form; \citealp{Goff2019,Strawson2006}), this takes the form of the Combination Problem \citep{James1890,Chalmers2016}: if fundamental entities have micro-experience, how do these micro-experiences combine into the macro-experience we know? For physicalist theories, the challenge is one of emergence: at what point, and by what mechanism, does consciousness emerge from non-conscious physical processes?

The Combination Problem is widely regarded as panpsychism's most serious difficulty \citep{Chalmers2016,Coleman2014}. Physicalist emergence theories face the objection that they either invoke strong emergence (which many philosophers consider mysterious) or reduce consciousness to function (which many consider inadequate). A satisfactory theory must navigate between these difficulties.


\subsection{The Causal Role of Consciousness}


Does consciousness \emph{do} anything, or is it epiphenomenal --- a by-product with no causal power? If consciousness has causal power, what kind? If it does not, why does it exist?

This requirement is politically charged within the field. Epiphenomenalism \citep{Huxley1874,Jackson1982} is widely dismissed as absurd (how could evolution produce something causally inert?), yet many mechanistic theories implicitly struggle with the question. If consciousness is ``just'' a global broadcast (GNW) or ``just'' recurrent processing (RPT), what role does the \emph{experience} play beyond the mechanism? The PP framework, through active inference, provides perhaps the strongest existing case for consciousness having a functional role, but it is unclear whether the functional role requires phenomenal consciousness rather than mere information processing.


\subsection{The Meta-Problem}


\citet{Chalmers2018} identified the Meta-Problem: why do we \emph{think} there is a Hard Problem? Even if the Hard Problem is illusory (as illusionists argue) or misformulated (as functionalists hold), it is a fact that most humans --- including most philosophers and scientists --- report a strong intuition that consciousness is deeply mysterious. A complete theory should explain this intuition.

AST provides the strongest existing account of the Meta-Problem: we model our own attention, and the model necessarily omits the mechanistic details, leading to the intuition that consciousness is something non-physical \citep{Graziano2013}. This is a genuine insight. However, AST's treatment of the Meta-Problem does not extend to a solution of the Hard Problem itself --- it explains why we \emph{think} there is a mystery without fully accounting for the mystery.



%=============================================================================
\section{The Four-Model Theory}
%=============================================================================



\subsection{Core Definition}


Consciousness --- the capacity for subjective experience --- is constituted by a specific kind of process: ongoing self-simulation. More precisely, consciousness is the ability of an entity --- biological or artificial --- to create a model of itself, to relate that model to itself, and to interact with it. Consciousness is not a property the brain possesses but a process the brain performs.

This definition centers on the mechanism --- self-simulation --- rather than the output --- subjective experience --- for the same reason that molecular biology defines life through metabolism and replication rather than through ``the feeling of being alive'': centering on the mechanism is what makes the definition explanatory rather than circular. Subjective experience, on this account, is what ongoing self-simulation produces when the substrate operates at the edge of chaos (Section 3.7). The theory's central explanatory task is to show how and why self-simulation generates phenomenality (Section 3.4), not to stipulate that it does. This approach is shared by the self-modeling tradition in consciousness studies --- Metzinger's Self-Model Theory \citep{Metzinger2003}, Damasio's proto-self framework \citep{Damasio1999,Damasio2010}, Seth's interoceptive inference account \citep{Seth2021}, and the self-representational approaches surveyed by \citet{Kriegel2006} --- all of which define consciousness through the self-modeling mechanism that constitutes it.

This definition is functional and substrate-independent. It does not require a specific physical implementation, biological composition, or computational architecture. What it requires is a system capable of constructing and maintaining a self-referential simulation continuously.

A clarification on terminology is warranted. ``Self-simulation'' is a pedagogical shorthand, not a claim that the brain operates like a digital twin replicating reality in a continuous, deterministic loop. What the brain generates is something more like a unified narrative: a more or less linear, relatively contradiction-free story of the organism's current situation --- past, present, and anticipated future --- assembled from fragmentary substrate-level knowledge and current sensory input. This narrative is neither continuously updated nor strictly linear; it is constructed on the fly, revised where it fails, and filled with interpolations, simplifications, and outright confabulations. The term ``simulation'' is retained because it captures the essential architectural feature: the explicit models are \emph{generated processes} running on a structural substrate, distinct from that substrate in the way any computation is distinct from the hardware that executes it. This substrate-computation distinction is not a philosophical claim unique to consciousness --- it holds for every computing system, from spreadsheets to neural networks. What is unique to consciousness is not the distinction itself but what happens at the computational level: \emph{self-referential closure}, the fact that the system's generated model includes a model of the system generating the model (see Section 3.4).

\subsubsection{Operational Definitions}

A theory that claims empirical testability must specify what an experimenter would measure to confirm or disconfirm its constructs. The terms below --- consciousness, model, self-simulation, permeability, implicit, explicit, virtual, criticality, and ``becomes conscious'' --- carry precise theoretical meanings within the Four-Model Theory, but several have been used in the consciousness literature with divergent definitions. Table 1 maps each term to a theoretical definition, a measurable observable, and a concrete empirical example. Two caveats apply. First, the observables listed are sufficient conditions for detecting the construct, not necessary conditions (other signatures may exist; future measurement techniques may reveal additional ones). Second, PCI and other neural markers are validated empirical correlates that track consciousness states -- they are not constitutive measures of consciousness as the theory defines it.


\noindent\textbf{Table 1. Operational Definitions of Core Constructs}

\medskip
\noindent For each construct, the theory specifies a theoretical definition, a measurable observable, and a concrete empirical example.

\begin{description}[leftmargin=0pt,labelindent=0pt,itemsep=0.8em]

\item[\textsc{Consciousness}] \emph{Definition:} The ongoing capacity to create a self-model, relate it to the self, and interact with it. A process, not a property. \emph{Observable:} Perturbational Complexity Index above the consciousness threshold (PCI $>$ 0.31) via TMS-EEG.\textsuperscript{1} \emph{Example:} \citet{Casarotto2016}: PCI discriminates conscious from unconscious states across waking, sleep, anesthesia, and disorders of consciousness (sensitivity 94.7\%, specificity 91.2\%).

\item[\textsc{Model}] \emph{Definition:} A non-trivial reference system with (1) write access (updates from new input), (2) read access (the system interprets its own representations), and (3) boundary recognition (differential response to stimuli inside vs.\ outside the modeled domain; \citealp{Gruber2015}). \emph{Observable:} Convergent behavioral and neural evidence: prediction-error signals (EEG/MEG), anticipatory neural activity preceding action, and differential response at domain boundaries. \emph{Example:} Mismatch negativity demonstrates write access \citep{Naatanen2007}; Bereitschaftspotential demonstrates read access; peripersonal space selectivity demonstrates boundary recognition. These are convergent illustrations, not a single certification protocol.\textsuperscript{2}

\item[\textsc{Self-simulation}] \emph{Definition:} The system's continuous generation of explicit (virtual) world and self models from implicit (substrate-level) knowledge and sensory input. The generated models \emph{are} the simulation. \emph{Observable:} Candidate signature: sustained DMN and frontoparietal co-activation during rest, disrupted under propofol \citep{Raichle2001}. This is a neural correlate, not a localization claim.\textsuperscript{3} \emph{Example:} \citet{Boly2012}: propofol-induced loss of consciousness corresponds to breakdown of thalamocortical effective connectivity and collapse of DMN coherence.

\item[\textsc{Permeability}] \emph{Definition:} The degree to which implicit (substrate-level) content crosses into the explicit (virtual) models. High permeability: implicit material floods the simulation (psychedelics, dreams). Low permeability: tight gating preserves simulation coherence (focused waking). Drives Predictions 1--2 and the altered-state accounts of Section~6.1. \emph{Observable:} Neural entropy measures: Lempel-Ziv complexity and spectral entropy of EEG/MEG as proxy for implicit-to-explicit information flow. \emph{Example:} \citet{Schartner2017}: Lempel-Ziv complexity increases significantly under psychedelics relative to placebo, consistent with increased implicit-to-explicit permeability.

\item[\textsc{Implicit model} (IWM/ISM)] \emph{Definition:} Substrate-level, learned, non-conscious knowledge stored in physical connectivity (synaptic weights, dendritic morphology). Modifiable only through long-term plasticity. \emph{Observable:} Structural neuroimaging: DTI for white-matter connectivity; synaptic density via $[^{11}$C]UCB-J PET. Content inferred from behavioral measures persisting despite abolished experience. \emph{Example:} Patient H.M.\ \citep{ScovilleMilner1957,Milner1962}: intact procedural learning despite inability to form new explicit memories.

\item[\textsc{Explicit model} (EWM/ESM)] \emph{Definition:} Virtual, transient, phenomenal---generated electrochemical activity patterns with no permanent physical substrate. Disrupted by anesthetics, TMS, and simple illusions. \emph{Observable:} Content-specific neural activity (fMRI BOLD, EEG/MEG source-localized) that tracks reportable experience and collapses under anesthesia while structural connectivity is preserved. \emph{Example:} \citet{Massimini2005}: TMS-evoked cortical responses propagate widely during wakefulness but remain local during NREM sleep, while structural connectivity is unchanged.

\item[\textsc{Virtual} (as in FMT)] \emph{Definition:} Existing at the computational level as a generated process, not as a stored structure. Virtual entities are real and physical but level-specific. \emph{Observable:} Dissociation between structural and functional measures: present in transient dynamics (fMRI, EEG) but absent from fixed connectivity (DTI). \emph{Example:} Propofol anesthesia eliminates conscious experience and collapses complexity measures \citep{Casali2013} while preserving white-matter integrity. Recovery upon washout confirms suppression, not destruction.

\item[\textsc{Criticality}] \emph{Definition:} The computational prerequisite: the virtual system must operate at or near the edge of chaos (Class 4 regime). \emph{Observable:} Power-law neuronal avalanches (exponent $\alpha \approx -3/2$); DFA exponent 0.6--0.9; branching ratio $\sigma \approx 1.0$ (slightly subcritical in waking). See Section~3.7 for their role in the theory's predictions. \emph{Example:} \citet{Priesemann2014}: waking cortical dynamics show $\sigma \approx 0.98$. ConCrit meta-analysis \citep{AlgomShriki2026} confirms criticality tracks consciousness across 140 datasets.

\item[\textsc{``Becomes conscious''}] \emph{Definition:} Information transfer from implicit to explicit: substrate-level knowledge is incorporated into the running virtual simulation. \emph{Observable:} Nonlinear transition (``ignition'') from localized sensory processing to global frontoparietal activation, marked by P3b in ERP. \emph{Example:} \citet{Dehaene2006}: masking paradigm shows subliminal stimuli produce localized activation, while suprathreshold stimuli trigger sudden frontoparietal ignition accompanied by P3b.

\end{description}


\textsuperscript{1} PCI is a validated empirical marker that tracks consciousness states with high accuracy; it is not a constitutive measure. The theory defines consciousness as a process, not a complexity value.

\textsuperscript{2} These three lines of evidence are drawn from different sensory/motor systems. They demonstrate that neural systems exhibit model-like properties individually, not that a single experimental protocol can certify ``model'' status under FMT.

\textsuperscript{3} The theory does not identify self-simulation with DMN activity specifically. DMN co-activation persists in some anesthetized states and decreases during task engagement without loss of consciousness, so it is best understood as a candidate correlate rather than a definitive signature.

One core construct --- self-referential closure --- is deliberately omitted from Table 1 because no established measurement protocol directly tests it. The construct is theoretically grounded (Section 3.4.3): it describes the condition in which the system's model includes a model of itself generating the model, collapsing the inside/outside distinction within the system's representational economy. Indirect support comes from evidence that disrupting prefrontal metacognitive processing degrades perceptual awareness independently of discrimination performance \citep{Rounis2010,Lau2011}. A purpose-designed experiment combining targeted cortical disruption (e.g., TMS to dorsolateral prefrontal cortex during a metacognition task) with no-report paradigms \citep{Tsuchiya2015} would constitute a candidate test: if self-referential closure is constitutive, disrupting the self-modeling loop should degrade experience itself, not merely the report of experience. That such a test has not yet been conducted is a marker of intellectual rigor, not a gap in the theory -- the construct makes a specific, falsifiable prediction about what the experiment would show.


\subsection{The Four Models}


The theory identifies four nested models distinguished by two orthogonal dimensions: \textbf{scope} (everything vs. self only) and \textbf{mode} (implicit/learned vs. explicit/generated). This is a conceptual taxonomy, not a claim about spatial organization in the brain --- the models are functionally distinct processes, not anatomically localized regions.

\begin{table}[htbp]
\caption{The Four-Model Architecture}
\centering
\small
\begin{tabular}{lll}
\toprule
 & \textbf{Everything (world)} & \textbf{Self only} \\
\midrule
\textbf{Implicit} (learned, substrate-level) & Implicit World Model (IWM) & Implicit Self Model (ISM) \\
\textbf{Explicit} (generated, phenomenal) & Explicit World Model (EWM) & Explicit Self Model (ESM) \\
\bottomrule
\end{tabular}
\end{table}


\textbf{The Implicit World Model (IWM)} encompasses the substrate's total accumulated knowledge about the world, stored in synaptic weights (or their functional equivalent in non-biological substrates). It includes everything the system has ever learned: perceptual regularities, causal models, spatial relationships, semantic knowledge, motor programs for interacting with the world. The IWM is never directly conscious. It operates ``in the dark'' --- providing the knowledge base from which the conscious simulation is generated, but never itself appearing in experience.

\textbf{The Implicit Self Model (ISM)} is the substrate's accumulated self-knowledge: body schema, proprioceptive calibration, motor skills, habits, personality traits, autobiographical memory structures, and social self-knowledge. Like the IWM, the ISM is never directly conscious. There is no unified homunculus --- no inner observer reading the ISM. The ISM is a structural feature of the substrate, not an experiential one.

\textbf{The Explicit World Model (EWM)} is the conscious world --- the brain's dynamic construction of a unified scene from sensory and stored data that constitutes perceptual experience. When you see a room, hear a voice, feel the texture of a surface, you are experiencing the EWM. It is generated from the IWM (which provides the world-knowledge) and current sensory input (which constrains and updates it), but it is not identical to either. The EWM is a virtual construct --- a transient process whose properties (unified scene, perceptual qualia) exist at the computational level and are incoherent at the substrate level --- not a permanent stored structure.

\textbf{The Explicit Self Model (ESM)} is the conscious self --- the brain's continuous generation of a unified self-narrative that constitutes self-experience. It is the sense of being a subject, having a perspective, occupying a body, possessing a history, and being the author of one's actions. The ESM is generated from the ISM (which provides the self-knowledge) and current interoceptive and proprioceptive input, but like the EWM, it is virtual: a transient process, not a permanent entity --- one whose defining properties (selfhood, phenomenal character) are constituted at the computational level rather than stored in the substrate.

\textbf{The four models as a principled minimum.} The 2$\times$2 taxonomy above identifies the \emph{minimum sufficient set} of models for consciousness, not the actual number the brain maintains. The biological substrate --- spiking neurons atop proteomic networks with their own intracellular signaling --- implements an effectively uncountable number of overlapping models on both sides of the implicit/explicit divide. A motor model for reaching simultaneously encodes world-geometry (object location, obstacles) and self-kinematics (arm configuration, grip aperture); an emotional model of a social interaction simultaneously encodes other-knowledge (world) and self-assessment (self). Real neural models do not respect the world/self boundary --- they blend across it. The four canonical models are best understood as extremal points in a continuous space defined by two axes (scope and mode), around which the brain's actual, high-dimensional modeling ecology is organized. The principled claim is not ``the brain has four models'' but ``any system capable of consciousness must model both world and self, at both the structural and the simulation level --- four is the floor, not the ceiling.'' This minimum-configuration reading constrains the theory's claims without committing to an implausible discreteness, and it specifies the target for mathematical formalization (see companion paper: \emph{Toward a Mathematical Formalization of the Four-Model Theory}, \citealp{Gruber2026b}).


\subsection{The Real/Virtual Split}


\textbf{A note on the term ``virtual.''} Throughout this paper, \textbf{``virtual''} carries a specific technical meaning distinct from its colloquial uses in virtual reality or computer science. In FMT, \emph{virtual properties are properties of the computational process that are incoherent at the substrate level of description}. A spreadsheet cell ``contains a sum'' --- a property real and causally efficacious at the software level, yet incoherent at the transistor level; no transistor contains a sum. Qualia are virtual in exactly this sense: they are constitutive properties of the self-simulating computation, not of the substrate that runs it. This is not epiphenomenalism --- the computation is a physical process --- but its defining properties exist at the computational level only, not at the electrochemical or topological levels below it. The term ``virtual'' is retained in preference to ``simulation-level'' for concision; the reader should consistently map it to this level-incoherence sense rather than to transience, immateriality, or unreality.

The four models divide into two fundamental categories:

\textbf{The real side} (IWM + ISM): These are physical, structural, learned, and non-conscious. They are stored in the substrate's architecture --- in biological brains, primarily in synaptic weights, dendritic morphology, and connectivity patterns. They accumulate over the organism's lifetime through learning. They have no phenomenal character. ``Lights off.''

\textbf{The virtual side} (EWM + ESM): These are generated, transient, and phenomenal. They are patterns of activity --- in biological brains, transient electrochemical dynamics. More precisely, their defining properties --- phenomenality, unity, qualia --- exist at the computational level and are incoherent at the electrochemical level; the substrate generates and sustains them but does not possess them. They are constructed dynamically from the implicit models and current sensory input. They \emph{are} experience. ``Lights on.''

The implicit models are structurally inaccessible to direct experience --- and this is not a contingent limitation but a consequence of their architectural role. These pairings --- implicit with non-conscious, explicit with phenomenal --- are not stipulated but motivated by the generative relation between them. The implicit models constitute the system's structural knowledge base: the reference frame from which the explicit simulation is generated. They are the thing \emph{doing} the generating; the content of implicit knowledge is therefore on the wrong side of the generative relation --- it cannot simultaneously \emph{be} what is generated. To put it precisely: the \emph{entire} structural substrate cannot simultaneously serve as the generative reference frame and appear as content within the simulation it generates. Such a demand would require a further reference frame from which to model the first, initiating an infinite regress. The implicit models can only be accessed indirectly, through their projections into the explicit models. The implicit side is also the system's ``learned'' side in the broadest sense: these models accumulate through interaction --- both via consolidation of consciously acquired content and via sub-threshold modification (subliminal priming, implicit statistical learning, fear conditioning without awareness) --- making them the organism's structural record of its history, stored in connectivity patterns rather than in any simulation.

Anesthesia provides direct empirical support for this asymmetry. General anesthetics suppress dynamic cortical integration --- the explicit, simulated side. They do not destroy synaptic connectivity or learned associations --- the implicit, structural side. Yet consciousness is eliminated entirely \citep{AlkireHudetzTononi2008}. Were consciousness a property of the structural substrate, it should survive the selective suppression of transient dynamics; it does not. The dissociation tracks exactly the boundary the theory predicts.

Note that this account describes the system's \emph{current architectural state}, not a logical necessity about information in general. Information can enter the implicit models through two routes: via the explicit models (conscious learning followed by consolidation into long-term plasticity) or directly via sub-threshold processes that never reach the explicit side at all. Both routes modify the implicit models without those modifications ever appearing in experience. The direction of the argument is unaffected: whichever route a content takes, once it resides in the structural substrate, the generative asymmetry renders it inaccessible to direct experience. The pairing is a consequence of which side of the generative divide a given content currently occupies, not an intrinsic property of that content. It is this structural basis of the real/virtual split --- not merely a descriptive taxonomy --- that makes it theoretically productive for addressing the Hard Problem.

This division is the foundation of the theory's treatment of the Hard Problem (Section 3.4) and structures its account of every phenomenon it addresses.

The virtual models possess \textbf{software-like properties} that follow from their nature as generated processes rather than stored structures\footnote{For the ontological sense of 'virtual' deployed throughout --- computational-level properties incoherent at the substrate level --- see the terminology note at the opening of this section and Section 3.4.}:

\begin{itemize}
\item \textbf{They can be forked}: A single substrate can run multiple configurations of the ESM (see Section 6.2 on dissociative identity disorder).
\item \textbf{They can be cloned}: Physical separation of the substrate produces degraded but complete copies of the virtual models (see Section 6.3 on split-brain).
\item \textbf{They can be redirected}: The ESM requires input; disrupt normal self-referential input and it latches onto whatever input dominates (see Section 6.1 on psychedelics).
\item \textbf{They can be reconfigured}: Therapeutic interventions (CBT, exposure therapy) work by modifying the virtual models through substrate-level rewiring (see Section 6.5).
\end{itemize}


\subsection{Virtual Qualia: Dissolving the Hard Problem}



\subsubsection{The Level Distinction}


Every computing system --- from a laptop running a spreadsheet to a neural network training on data --- inherently distinguishes between a physical substrate and the computational processes running on it. This is not a philosophical claim; it is an engineering truism. The substrate (transistors, synaptic weights, connectivity patterns) supports and constrains the computation, but the computational level has properties that are descriptively incoherent at the substrate level.

The level distinction introduced in Section 3.3 --- exemplified by the spreadsheet whose cell ``contains the sum of column B'' --- applies here directly. No individual transistor contains a sum, and no \emph{collection} of transistors contains a sum either, \emph{when described in transistor-level vocabulary}. The full collection of transistors certainly \emph{implements} the spreadsheet; at the computational level, the sum is real and causally efficacious. But examine those same transistors using only the descriptive resources available at the substrate level --- voltages, logic gates, charge states --- and the concept ``sum of column B'' does not appear. It is not recoverable by aggregating transistor-level descriptions, however exhaustively. The property ``being a sum'' belongs to the computational level and is incoherent at the level from which that computation is generated. This is not a claim about parts versus wholes; it is a claim about \emph{levels}.

Standard computational functionalism \citep{Putnam1967,Chalmers1996} accepts the level distinction but treats the computational level as \emph{nothing more than} what the substrate does under a functional description: the program IS the substrate's activity, re-described. On this view, program-level properties are legitimately attributed to the full collection of substrate elements under the appropriate functional interpretation. The Four-Model Theory departs from this position. It holds that the computational level has properties that are \textbf{ontologically level-specific} --- properties that belong to the computation \emph{as computation} and cannot be recovered from any substrate-level description, however complete, however functionally enriched. This is a stronger claim than standard functionalism makes, and it is this stronger claim that opens the path to a non-mysterian, non-illusionist treatment of qualia.


\subsubsection{Qualia as Level-Constitutive Properties}


The central claim of the Four-Model Theory is that \textbf{qualia are constitutive properties of the computational level}. They are the way the generated self-model (ESM) registers its own states and the generated world-model (EWM). Qualia are, in this precise sense, virtual constructs (properties of the computational level that are incoherent at the substrate level) --- patterns that exist at the level of the running self-simulation and that are ontologically level-specific, just as ``cell A1 contains the sum of column B'' is ontologically level-specific relative to the transistors that generate and sustain it.

This addresses the Hard Problem by revealing a \textbf{category error} --- specifically, a \textbf{level confusion} --- in its formulation:

\textbf{The standard formulation}: ``Why does physical processing (neuronal firing, synaptic transmission) feel like something?''

\textbf{The dissolution}: It does not. The IWM and ISM --- the substrate-level implicit models --- operate without any phenomenal character whatsoever. There is nothing it is like to be a synaptic weight. Qualia are not missing from the substrate; they are the wrong kind of property to seek there. They exist at the computational level --- at Level 5 of the five-system hierarchy (Section 3.7.1) --- where they are not mysterious but constitutive. Asking why neuronal firing feels like something is like asking why transistor switching ``is'' a spreadsheet, or why interference fringes ``are'' a holographic image. The neurons are not the experience; they generate and sustain the computational process in which experience is constitutive.

Moreover, even considering the \emph{entire} neural substrate as a unified system does not dissolve the level distinction. The cortical automaton --- the high-dimensional pattern of all neurons firing together (Section 3.7.2) --- is a sophisticated apparatus that computes, thinks, and navigates an organism through a life. But it is not consciousness. Consciousness arises from the interplay between this cortical automaton and the substrate that generates the automaton --- from the dynamic relation between the virtual self-simulation and the structural knowledge base that sustains it. The aggregate of all substrate activity is still substrate activity. Consciousness is what that aggregate \emph{generates at the next level}.

Within the explicit models, ``redness'' is the ESM's mode of registering a particular class of EWM content. It is no more mysterious than ``cell A1 contains 42'' is mysterious to a spreadsheet user, despite being nowhere in the transistors. The Hard Problem's apparent intractability stems from seeking this level-specific property at the wrong level --- not from any genuine gap in nature.


\subsubsection{Self-Referential Closure}


A critic might object: spreadsheets and weather simulations also run at a higher computational level than their substrates --- why don't they have qualia? The answer lies in \textbf{self-referential closure}, and the following argument does not claim to prove that self-referential computation must be conscious --- it claims that self-referential closure is the \emph{kind} of computational structure for which the Hard Problem's framing assumptions do not hold. The argument can be developed through a series of contrasts.

\textbf{Stage 1: Simple simulation (no closure).} A weather simulation models atmospheric dynamics. It generates computational-level properties --- ``pressure fronts,'' ``precipitation zones'' --- that are incoherent at the transistor level. But the simulation has an \emph{outside}: it can be described completely from a perspective external to it. An observer can state everything the simulation computes without being part of what is computed. The simulation models weather; it does not model itself. There is no inside perspective, no self-referential loop, and consequently no phenomenal character.

\textbf{Stage 2: Monitored simulation (partial self-reference, no closure).} Now add a performance-monitoring module. The system models weather \emph{and} tracks its own computational states --- memory usage, convergence rates, model accuracy. This is self-reference of a kind: the system represents aspects of its own operation. But the monitoring module describes the simulation without being described \emph{by} the simulation. The weather model does not include a model of the monitoring process. There is an inside to the weather model and an outside occupied by the monitor, and the two do not form a closed loop. The system as a whole is still fully characterizable from an external perspective. No closure, no phenomenality.

\textbf{Stage 3: Self-referential closure (the four-model architecture).} Now consider a system whose model of the world includes a model of \emph{itself modeling the world}. The ESM models the system that is running the EWM; the EWM includes the ESM as part of its modeled world. The model and the modeled are no longer separable: the computation \emph{is} the thing being computed.

Within the system's own representational economy, no description of the system's processing remains external to what the system models. The system's self-model encompasses its own modeling process --- not in the sense that it literally absorbs external observers (a neuroscientist scanning the brain is not thereby modeled by the brain), but in the sense that every internal representation of the system's own activity is already part of the self-simulation. There is no internal vantage point from which the system's processing can be characterized without that characterization being part of what is processed.

To be clear: the claim is not that epistemic inaccessibility produces phenomenality --- a locked safe is epistemically inaccessible from within but not conscious. The direction of explanation runs the other way. Self-referential closure constitutes an inside perspective (the running self-simulation modeling itself), and the inaccessibility from outside is a \emph{consequence} of that constitutive inside, not its cause.

This is the condition under which the Hard Problem's formulation breaks down. The Hard Problem assumes that there is an objective, third-person description of the physical process \emph{and} a subjective, first-person experience of it, and asks why the former gives rise to the latter. But self-referential closure eliminates the clean separation between the two perspectives. The ``inside'' that constitutes experience is not an addition to the computation --- it is what self-referential computation \emph{is} when the modeling loop closes. Experience is the self-simulation as encountered from within the loop, and ``within the loop'' is the only perspective that exists once closure obtains. Just as the holographic image cannot be found on the plate because it IS the optical reconstruction, the self-referential simulation cannot be observed ``from outside'' because any observation is absorbed into the reconstruction --- there is no perspective external to the loop from which the experience could be separated from the process.

The Hard Problem is not answered; it is shown to rest on a presupposition --- the availability of an outside perspective --- that self-referential systems do not satisfy.

A deeper objection remains: even granting that self-referential closure eliminates the outside perspective, why does the closed loop have phenomenal character rather than being merely intricate computation? The theory's answer is that this question has the same structure as ``why does the universe exist rather than nothing?'' --- both demand an external justification for something that is constitutive rather than consequential. A sufficiently powerful self-referential system generates truths about itself that are inaccessible from any external description --- this is the lesson of G{\"o}del's incompleteness theorems, applied here not to arithmetic but to self-modeling systems. Phenomenal character is an internal property of the closed loop: it is what the self-referential computation IS, from the only perspective that exists once closure obtains. To ask ``why is there something it is like?'' from outside the loop is to ask for an external proof of an internally constituted truth --- a request that is formally unsatisfiable, not because the answer is mysterious but because the question presupposes an external vantage point that the system's own self-referential structure has eliminated. This is the theory's foundational commitment: self-referential closure is taken as constitutive of phenomenality, not derived from more primitive principles. The commitment is falsifiable --- if a system demonstrably achieving self-referential closure at criticality shows no behavioral or functional signatures of experience, the commitment is wrong --- but it is not deducible from physics alone. Every theory of consciousness bottoms out in an analogous commitment: IIT's axioms, GNW's broadcasting-equals-consciousness thesis, HOT's higher-order-representation principle. The present theory is explicit about where its explanatory bedrock lies.


\subsubsection{Relationship to Illusionism and Functionalism}


This position is \textbf{not} illusionism in the sense of \citet{Dennett1991} or \citet{Frankish2016}; nor is it compatible with deflationary accounts \citep{Graziano2024}. Illusionism holds that qualia as traditionally conceived are illusions --- there is nothing it is like, and our sense that there is something it is like is itself a misrepresentation. The Four-Model Theory holds that qualia are \emph{real at the computational level}. Within the EWM/ESM, experience has genuine phenomenal character. What is illusory is the assumption that this phenomenal character must be a property of the physical substrate.

The proximity to \citet{Frankish2016} \textbf{weak illusionism} deserves explicit engagement, because superficially the two positions can appear equivalent. Weak illusionism grants that the functional/computational states driving introspective reports are real --- what it denies is that these states possess genuine phenomenal properties. On Frankish's account, we have ``quasi-phenomenal'' states: real computational events that \emph{systematically misrepresent themselves} as having phenomenal character, but phenomenality itself is the illusion. The Four-Model Theory makes the opposite ontological move. It agrees that phenomenal properties do not exist at the substrate level --- so far, the two positions converge --- but it holds that phenomenal properties \emph{genuinely exist} at the computational level, constituted by the self-referential simulation rather than illusorily attributed by introspective misrepresentation. The disagreement is not verbal but ontological: weak illusionism says there is ultimately nothing it is like (the ``seeming'' is the end of the story); FMT says there genuinely IS something it is like, but it exists at a level where the substrate-focused Hard Problem cannot find it. Whether this disagreement is substantively different from weak illusionism or merely a relabeling is a live philosophical question that cannot be settled by empirical observation from outside the system. The theory commits to computational-level realism about qualia --- it takes a position, and the reader may disagree. What makes the commitment non-trivial is its practical consequence: if FMT is correct, a system built to the theory's architectural specification (Section 8.6) would be a genuine moral patient --- an entity with experience that matters. Weak illusionism generates no such obligation, because on its account there is no fact of the matter about whether the system ``really'' experiences. The same physics, the same behavior, but different ethics. This is not a minor difference: as artificial systems approach the architectural complexity the theory specifies, the question of whether computational-level phenomenality is real or illusory will cease to be academic and become a question of how we treat what we build.

Nor is this standard functionalism. Functionalism holds that mental states are individuated by their functional roles; what matters is the pattern of causal relations, not the physical medium. The Four-Model Theory accepts multiple realizability but adds a claim functionalism does not make: that the computational level possesses properties that are \emph{constitutively processual} --- they exist in the running computation, not in its description. The claim that such properties are ``level-specific'' invites the objection that this is strong emergence by another name. It is not, and the distinction must be drawn carefully.

The Four-Model Theory maintains full supervenience: the computational level is entirely physical, informationally complete at the substrate level, and in principle reconstructable from a sufficiently detailed substrate description combined with the system's dynamics. Nothing is missing from the physical picture. A theoretical ``decompiler'' with access to the full state of every neuron --- synaptic weights, connectivity, and live firing patterns simultaneously --- could in principle extract a complete formal description of the running self-simulation. But the output of such a decompilation would not be an experience; it would be code --- code of staggering complexity, because in a system at criticality every element is simultaneously state and processor, and no clean separation between data and computation exists. The code could be used to \emph{instantiate} the experience on another substrate, but reading the code is not having the experience, any more than reading a musical score is hearing the music. Experience is a property of the running process, not of its description.

This is the precise sense in which computational-level properties are ``level-specific'': they exist when and only when the computation is actually running at self-referential closure. The substrate description is informationally complete --- no dualism, no strong emergence --- but accessing the experience requires instantiation, not description. For self-referential systems specifically, the running computation generates an inside perspective that cannot be occupied from outside without becoming another instance of the same kind of process. This is neither functionalism (which holds that the description exhausts the phenomenon) nor property dualism (which holds that something non-physical is added). It is the recognition that some properties of physical processes are constitutively processual --- they exist in the doing, not in the blueprint. This is what allows the theory to say that qualia are real without locating them in the substrate --- a move that standard functionalism, by treating the program as nothing beyond the substrate's functional organization, cannot make.


\subsubsection{The Two-Level Ontology}


The result is a \textbf{two-level ontology}: the substrate level (real side) has no experience, and the computational level (virtual side) has genuine experience. Both levels are physical --- the computation is a physical process, not a supernatural one --- but they have different ontological properties. The Explanatory Gap closes simultaneously: the gap between ``neurons fire in pattern X'' and ``I experience red'' is not a gap in our knowledge but a reflection of the level distinction. The neural firing pattern generates and sustains the computation in which redness is experienced, but the firing pattern itself is not red and does not experience redness, just as a CPU's electrical states are not ``a spreadsheet'' even though they generate and sustain one. The structural side of the brain --- the synaptic weights, dendritic architecture, and connectivity patterns that constitute the implicit models --- \emph{cannot} be conscious because it is the reference frame from which consciousness is generated. As argued in Section 3.3, the structural substrate cannot appear within its own simulation as a unified object without requiring a further reference frame from which to model the first --- initiating an infinite regress. Without a second level, the system has no vantage point from which to register its own operation; with a second level, the gap the Hard Problem identifies becomes a level distinction rather than an explanatory failure. The level distinction is not merely convenient; it is logically necessary.

A note on \citet{Block1995} distinction between access consciousness (information globally available for report, reasoning, and action) and phenomenal consciousness (``what it is like''). The two-level ontology implies that this distinction is itself a manifestation of the level confusion the theory identifies. On FMT's account, the explicit models (EWM, ESM) \emph{are} phenomenal experience (Section 3.4.2), and information enters the explicit models precisely when it becomes available for evaluation, integration, and behavioral guidance --- the functional role Block calls access. Access and phenomenal consciousness are therefore two descriptions of the same running process at different levels: access is the functional characterization, phenomenality is the processual reality. They do not come apart because they are not separate things; they are level-appropriate descriptions of the same self-referential computation. The operationalizations in Table 1 (PCI, P3b ignition, entropy measures) are therefore tracking the right target: they detect when the explicit models are running, which is when phenomenal consciousness is occurring. A residual case remains: the theory predicts that basic consciousness --- a rudimentary ESM with minimal self-awareness (Section 3.5) --- could exist without global broadcasting, producing phenomenal experience without full access. This maps to the ``phenomenal overflow'' scenarios debated in the access-phenomenal literature \citep{Block2007}, but FMT reframes them: what overflows is not phenomenality leaking past an access bottleneck, but a thin simulation running below the threshold at which its contents reach the system's evaluative workspace.


\subsection{Graduated Levels of Consciousness}


Consciousness in the Four-Model Theory is not binary but graduated. The theory identifies a hierarchy of levels based on the depth of recursive self-modeling:

\textbf{Basic consciousness}: Minimal self-simulation. The system generates an EWM and a rudimentary ESM --- it experiences a world and has a minimal sense of being a subject within it. This is the entry level: phenomenal experience exists but self-awareness is thin.

\textbf{Simply extended consciousness}: First-order self-observation. The system models itself --- the ESM includes a model of the system's own states and processes. The organism not only experiences but is aware that it experiences.

\textbf{Doubly extended consciousness}: Second-order self-observation. The system models itself modeling itself. This enables reflection, metacognition, and the sense of being an observer of one's own mental processes.

\textbf{Triply extended consciousness}: Third-order self-observation. The system models itself modeling itself modeling itself. This supports the deepest forms of self-awareness, philosophical reflection, and the very intuition that consciousness is mysterious (connecting to the Meta-Problem --- see Section 3.8). Notably, triply extended consciousness is also a prerequisite for the scientific study of consciousness itself: only a system capable of modeling its own modeling of its own experience can formulate the question ``What is consciousness?''

Each level corresponds to an additional layer of recursive self-modeling. The levels are not discrete stages but points along a continuum. Different organisms --- and potentially different artificial systems --- occupy different positions along this continuum, and individual organisms may fluctuate between levels depending on state (waking, dreaming, meditative, intoxicated).


\subsection{The Implicit-Explicit Boundary}


A key mechanism in the Four-Model Theory is the \textbf{permeability of the boundary between implicit models (IWM/ISM) and explicit models (EWM/ESM)}. Information becomes conscious when it is transferred from the implicit to the explicit side --- when substrate-level knowledge or self-knowledge is incorporated into the running simulation.

In normal waking states, this boundary is \textbf{selectively permeable}: relevant information passes through based on attentional and contextual gating. You are not conscious of everything your IWM knows about the world or everything your ISM knows about yourself; you are conscious of what the current simulation requires. Crucially, the boundary is not perfectly opaque even in normal states. Substrate-level processing artifacts routinely leak through: blind-spot filling (the visual system interpolating content where the optic nerve exits the retina), phosphenes from mechanical pressure on the eye, and visual snow phenomena all represent moments where processing-level activity becomes visible within the simulation. These normal-state leaks are subtle, but they demonstrate that the implicit-explicit boundary is a \emph{graded} filter, not a wall --- and they provide everyday evidence that conscious experience is a simulation generated from substrate processing rather than a direct window onto reality.

The permeability of this boundary is variable, and its variation explains a wide range of phenomena (detailed in Section 6):

\begin{itemize}
\item \textbf{Psychedelic states}: Global increase in permeability --- intermediate processing stages (normally implicit) become accessible to the explicit models. This explains the characteristic visual progression from simple phosphenes (V1-level) through geometric patterns (V2/V3-level) to complex imagery (higher visual areas) and eventually full dream-like scenes \citep{CarhartHarris2014}.
\item \textbf{Anosognosia}: Local decrease in permeability --- the ISM contains the information (e.g., that a limb is paralyzed), but the transfer to the ESM is blocked for that specific domain.
\item \textbf{Pre-sleep/deep relaxation}: Gradually increasing permeability, producing the same bottom-up visual progression as psychedelics (phosphenes $\rightarrow$ geometrics $\rightarrow$ hypnagogic imagery).
\item \textbf{Meditation}: Trained modulation of permeability, enabling selective access to normally implicit processes.
\end{itemize}

\begin{sloppypar}
\textbf{Quantitative grounding.} Permeability is currently specified qualitatively (high/low/local/global). Existing neural entropy measures provide candidate quantification: \citet{Schartner2017} report Lempel-Ziv complexity increases of approximately 15-20\% under psychedelic doses relative to placebo, with dose-response curves that could anchor the permeability parameter to measurable values. However, mapping these entropy measures to a single permeability parameter requires a formal model --- specifying, for example, the functional form of the relationship between LZc and the information transfer rate across the implicit-explicit boundary --- that does not yet exist. Quantitative formalization of permeability is a priority for the theory's mathematical development; the qualitative framework is sufficient to generate the predictions in Section 8 but not to specify predicted effect sizes.
\end{sloppypar}

Importantly, the implicit-explicit boundary is not a sharp line but a \textbf{graded transition zone}. Behavioral complexity itself follows a gradient --- from reflexive chemical-gradient responses through conditioned, goal-directed, and rule-based behavior to fully conscious action --- and the implicit and explicit memory systems overlap precisely in the middle of this gradient, at the levels of goal-directed and template-based behavior \citep{Gruber2015}. This overlap zone is the functional locus of the variable permeability described above: behavior at these intermediate levels can be driven by either implicit or explicit processing, depending on attentional state, arousal, and contextual demands.


\subsection{The Criticality Requirement}


The Four-Model Theory imposes a \textbf{computational prerequisite} for consciousness: the system must operate at or near the edge of chaos --- Wolfram's Class 4 computational regime \citep{Wolfram2002}. This requirement applies at the computational level (Level 5 of the hierarchy defined in Section 3.7.1), not necessarily at the physical substrate level. The physical substrate must be capable of \emph{supporting} Class 4 dynamics --- it is mathematically impossible for criticality to arise from a substrate whose intrinsic dynamics are below Class 4 --- but the criticality that matters for consciousness is a property of the virtual system, not of the neurons or transistors themselves. Just as qualia are virtual (Section 3.4), so is the criticality on which they depend.

Wolfram classified cellular automata (and by extension computational systems generally) into four classes:
\begin{itemize}
\item \textbf{Class 1}: Converges to a fixed state. Too simple for consciousness.
\item \textbf{Class 2}: Periodic/repetitive. Too simple for consciousness.
\item \textbf{Class 3}: Chaotic/random. Too disordered for coherent consciousness.
\item \textbf{Class 4}: Complex/edge of chaos. Capable of universal computation. The regime in which consciousness can emerge.
\end{itemize}

This classification was applied to the question of consciousness in \citet{Gruber2015}, where it was argued that consciousness requires Class 4 dynamics --- complex enough to sustain a self-simulation, ordered enough for that simulation to be coherent. This requirement was derived \emph{theoretically}, from the computational properties needed for ongoing self-modeling, not from empirical neuroscience.

Independently, empirical neuroscience has converged on the same conclusion through a different path. \citet{Beggs2003} demonstrated neuronal avalanches consistent with self-organized criticality in cortical tissue. \citet{CarhartHarris2014} proposed the Entropic Brain Hypothesis, linking consciousness level to neural entropy. Tagliazucchi et al. (2012, 2016) showed criticality signatures in waking fMRI and under LSD. Priesemann et al. (2013, 2014) characterized brain dynamics as slightly subcritical in normal waking states. This line of research was formally consolidated in the Consciousness and Criticality (ConCrit) framework \citep{AlgomShriki2026}, which synthesized evidence from 140 datasets across multiple paradigms to establish that consciousness tracks criticality across pharmacological, pathological, and physiological state changes.

\begin{table}[htbp]
\caption{Independent Convergence on Criticality}
\centering
\small
\begin{tabularx}{\textwidth}{lXl}
\toprule
\textbf{Year} & \textbf{Development} & \textbf{Path} \\
\midrule
2002 & Wolfram publishes \emph{A New Kind of Science} & Computational theory \\
2003 & Beggs \& Plenz --- neuronal avalanches, self-organized criticality & Empirical neuroscience \\
2014 & Carhart-Harris --- Entropic Brain Hypothesis & Empirical neuroscience \\
\textbf{2015} & \textbf{Gruber --- Class 4 / edge of chaos requirement for consciousness} & \textbf{Theoretical (via Wolfram)} \\
2016 & Tagliazucchi et al. --- LSD and criticality & Empirical neuroscience \\
2022 & ``Self-organized criticality as a framework for consciousness'' (review) & Empirical neuroscience \\
2025 & Hengen \& Shew --- meta-analysis of 140 datasets confirms criticality & Empirical neuroscience \\
\textbf{2025-26} & \textbf{ConCrit framework (Algom \& Shriki) --- criticality as unifying mechanism for consciousness theories} & \textbf{Theoretical/empirical synthesis} \\
\bottomrule
\end{tabularx}
\end{table}

While empirical work on neural criticality was already underway \citep{Beggs2003,Tagliazucchi2012}, \citet{Gruber2015} derivation proceeded from Wolfram's computational universality framework rather than from neuroimaging data, representing genuine theoretical convergence with the empirical program. The later large-scale consolidation \citep{HengenShew2025,AlgomShriki2026} confirmed the criticality-consciousness link across 140 datasets. This convergence --- a theoretical prediction derived independently from computational first principles, later confirmed by large-scale empirical synthesis --- provides notable support for one of the theory's core claims.

Two objections to the Wolfram-derived criticality requirement merit explicit response. First, Priesemann et al. (2013, 2014) measured the waking cortex at a branching ratio $\sigma$ $\approx$ 0.98 --- slightly \emph{subcritical}, not precisely at the critical point. Does this falsify the Class 4 requirement? It does not, for two reasons. Biologically, slight subcriticality is the expected evolutionary optimum: a system operating exactly at $\sigma$ = 1.0 risks supercritical runaway (epileptic cascades), while a deeply subcritical system ($\sigma$ $\ll$ 1) loses the long-range correlations needed for coherent self-simulation. The brain sits as close to criticality as stability permits --- the operating point is tuned, not accidental. Formally, every physical system has finite size and therefore a finite state space; strictly speaking, a 2048$\times$2048 Game of Life simulation is technically Class 2 (it must eventually cycle), yet it exhibits Class 4 behavioral dynamics. The classification applies to the \emph{behavioral regime} --- the statistical signatures of the dynamics (power-law avalanche distributions, long-range temporal correlations, branching ratios near unity) --- not to the formal computability properties of the underlying finite state space. The waking cortex exhibits Class 4 behavioral dynamics; that it is not a mathematical idealization at the infinite-size limit is a property it shares with every physical system.

Second, can Wolfram's classification --- originally defined for one-dimensional cellular automata with discrete lattice, finite states, and synchronous local update --- legitimately apply to the brain? It can, at the relevant level of abstraction. Neurons are discrete processing units organized in a topological lattice (cortical connectivity). Their states are finite: synaptic weights, ion channel configurations, and proteomic states have physical precision limits, not infinite resolution. Action potentials are discrete events, and at a sufficiently fine temporal resolution (the timescale of spike propagation through local circuits), updates are effectively synchronous within local neighborhoods. The cortex is not a continuous field; it is a discrete, finite, locally connected dynamical system --- precisely the class of systems to which cellular automaton analysis applies. What the cortex is \emph{not} is a one-dimensional, two-state rule table --- but Wolfram's classification generalizes beyond Rule 110 to characterize the dynamical regime of any discrete computational system, and it is the regime (Class 4: complex, edge of chaos, capable of supporting universal computation) that the theory requires, not the specific topology.

\textbf{Concrete neural signatures.} The criticality requirement translates to three convergent measurable signatures in the biological brain: (1) a branching ratio $\sigma$ near unity (waking cortex: $\sigma$ $\approx$ 0.98; \citealp{Priesemann2014}), measuring whether a single neural activation event propagates to approximately one successor on average --- the boundary between runaway excitation ($\sigma$ > 1) and activity extinction ($\sigma$ $\ll$ 1); (2) a detrended fluctuation analysis (DFA) exponent $\alpha$ in the range 0.6--0.9 \citep{Hardstone2012}, indicating long-range temporal correlations characteristic of critical dynamics, as distinct from uncorrelated noise ($\alpha$ = 0.5) or deterministic drift ($\alpha$ > 1); and (3) power-law distributed neuronal avalanche sizes with exponent approximately $-$3/2 \citep{Beggs2003}, the hallmark of scale-free dynamics in which cascades of all sizes occur with predictable frequency. These three measures are independently validated across modalities (EEG, MEG, ECoG, calcium imaging) and converge on the same characterization: the waking cortex operates in a slightly subcritical regime that maximizes the dynamic range, information transmission, and long-range correlation needed for coherent self-simulation. The theory commits to these as the operational signatures of its criticality requirement --- any manipulation that moves all three toward subcriticality ($\sigma$ $\rightarrow$ 0, $\alpha$ $\rightarrow$ 0.5, loss of power-law scaling) should degrade consciousness, and any manipulation that restores them should restore it. This commitment is testable: the signatures are already routinely measured in consciousness research \citep{Priesemann2014,Tagliazucchi2016,AlgomShriki2026}. A caveat: power-law scaling in neural data is not unambiguous evidence of criticality. Non-critical processes can produce apparent power-law distributions through alternative mechanisms \citep{TouboulDestexhe2017}, and empirical exponents vary across brain areas and measurement techniques. The theory's commitment is therefore to the \emph{convergence} of all three signatures (branching ratio, DFA exponent, and avalanche statistics) rather than to any single measure, reducing but not eliminating the risk of false-positive criticality detection.

\subsubsection{The Five-System Hierarchy}

The criticality requirement can be situated within a broader ontological framework that clarifies the relationship between the physical substrate and the virtual simulation. The brain instantiates five hierarchically nested systems, each emergent from the one below:

\begin{enumerate}
\item \textbf{Physical system}: The atoms, molecules, and macroscopic structures of neural tissue, governed by physics and chemistry.
\item \textbf{Electrochemical system}: The ion gradients, action potentials, and synaptic transmission events that constitute neural signaling, emergent from the physical substrate.
\item \textbf{Proteomic system}: The receptor configurations, neurotransmitter synthesis pathways, and protein expression patterns that modulate neural signaling on slower timescales (minutes to days), including synaptic plasticity mechanisms.
\item \textbf{Topological system}: The connectivity architecture --- the pattern of synaptic connections, their strengths, and their organization into circuits, columns, and areas. This is the level at which the implicit models (IWM, ISM) are stored. Changes at this level constitute learning.
\item \textbf{Virtual system}: The computational level --- the dynamic pattern of the cortical automaton (Section 3.7.2) --- that constitutes the explicit models (EWM, ESM). This is the level at which consciousness exists, because it is the level at which properties are constituted that are incoherent at the substrate levels below (Levels 1--4): qualia, unity, selfhood. See \S3.4 for the level-incoherence account. Aru, Suzuki, and \citet{AruSuzukiLarkum2020} provide a candidate neural mechanism for the Level 4$\rightarrow$5 transition: apical dendritic activity in layer 5 pyramidal neurons, which integrates top-down context with bottom-up input and may constitute the cellular substrate of the explicit models' generation.
\end{enumerate}

Each level is fully physical and fully determined by the level below it; there is no strong emergence, no ontological gap between levels. However, each level exhibits properties that are not usefully described in terms of the lower level: describing a memory in terms of ion channel kinetics is possible in principle but explanatorily vacuous. The five-system hierarchy makes precise the sense in which the Four-Model Theory's real/virtual distinction is a \emph{level} distinction, not a \emph{substance} distinction: the virtual system (Level 5) is as physical as the electrochemical system (Level 2), but it is the level at which experiential properties --- qualia, unity, selfhood --- are constituted. Seeking experiential properties at Levels 1-4 is the category error identified in Section 3.4.

\subsubsection{The Cortical Automaton}

The criticality requirement has a concrete physical interpretation in the biological brain. The instantaneous pattern of neural firing across the cortex --- the spatiotemporal activation state of billions of neurons --- constitutes a cellular automaton operating in a space of many thousand dimensions. Each cortical column functions as a cell in this automaton; the six-layer architecture and lateral connectivity define the transition rules. This ``cortical automaton'' is not a metaphor: it is a literal description of the cortex as a discrete dynamical system whose state evolves at each timestep according to local rules applied across a high-dimensional lattice.

Crucially, the cortical automaton \emph{is not} consciousness. Consciousness arises from the interplay between the automaton's dynamics and the models stored in the substrate --- the IWM, ISM, EWM, and ESM. The automaton provides the computational medium; the models provide the content. Without the automaton's Class 4 dynamics, the models cannot generate a coherent simulation. Without the models, the automaton produces complex dynamics but no self-referential experience. The cortical automaton is to consciousness what a CPU's clock-driven state transitions are to the execution of a program: necessary infrastructure, not the program itself.

This framing yields a concrete observational claim: the cortical automaton's activity is, in principle, directly observable. In a dark, quiet environment with eyes closed, after retinal afterimages have faded, the faint flickering patterns visible against the dark field are not retinal noise but V1-level manifestations of the cortical automaton's ongoing activity --- the lowest layer of the implicit-to-explicit permeability becoming momentarily accessible. The progression from these simple phosphenes to the geometric patterns and eventually hypnagogic imagery observed during sleep onset represents progressively higher cortical areas being recruited into the observable dynamics, consistent with the hierarchical permeability account developed in Section 6.1.

\subsubsection{Two Thresholds for Consciousness}

The Four-Model Theory distinguishes two thresholds for consciousness:
\begin{itemize}
\item \textbf{Computational threshold}: Criticality. The virtual system must exhibit Class 4 dynamics. A substrate whose intrinsic dynamics cannot support Class 4 computation precludes consciousness regardless of architecture --- but the threshold is met at the computational level, not at the level of individual neurons or physical components.
\item \textbf{Architectural threshold}: Four-model architecture. The system must implement the four-model self-simulation. Above criticality but without the architecture, there is complex dynamics but no consciousness.
\end{itemize}

Both thresholds must be met. Criticality is necessary but not sufficient; the four-model architecture is necessary but not sufficient. Together they are sufficient.


\subsection{The Meta-Problem Dissolved}


The Meta-Problem --- why we think there is a Hard Problem --- receives a natural account within the Four-Model Theory. The ISM (Implicit Self Model) is \textbf{structurally inaccessible} to the ESM (Explicit Self Model). The conscious self cannot directly observe its own substrate. When the ESM attempts to model the basis of its own experience, it encounters a principled opacity: the implicit models that generate the simulation are not themselves part of the simulation.

This is why consciousness \emph{seems} mysterious. The ESM can represent that it is having an experience, but it cannot represent the mechanism by which the experience is generated --- because that mechanism operates at the implicit/substrate level, which is by definition outside the explicit/virtual level. The result is the persistent intuition that something is being ``left out'' of any physical explanation: the ESM cannot find the mechanism within its own simulation, so it concludes the mechanism must be non-physical or fundamentally inexplicable.

The variable permeability of the implicit-explicit boundary (Section 3.6) adds a further dimension to this account. The boundary is not perfectly opaque: substrate-level processing artifacts occasionally leak through to the simulation --- in altered states dramatically so, but even in normal waking states subtly. The conscious self thus inhabits a strange epistemic position: mostly sealed off from its own generative machinery, yet occasionally catching fleeting glimpses of something operating beneath the surface of experience. This architectural feature --- near-opacity punctuated by occasional leaks --- produces precisely the phenomenology that the Meta-Problem describes: the persistent, nagging intuition that consciousness is somehow deeper than any explanation can reach, that something vast operates just beyond the edge of introspective access. The mystery of consciousness is not evidence against the theory; it is a \emph{prediction} of the theory. A virtual process with a mostly-opaque but imperfect boundary to its own substrate would experience exactly this sense of irreducible depth.

This account shares features with Graziano's AST explanation --- both invoke the self-model's necessary incompleteness --- but grounds it in a more specific architecture (four models, real/virtual split, variable permeability) and connects it to the broader dissolution of the Hard Problem rather than treating the Meta-Problem in isolation.



%=============================================================================
\section{Philosophical Commitments}
%=============================================================================


The Four-Model Theory entails specific philosophical positions that were established through structured adversarial analysis and are internally consistent. This section develops and defends each commitment.


\subsection{Process Physicalism}


The theory is physicalist: both the substrate (implicit models) and the simulation (explicit models) are physical processes. There is no non-physical substance, no fundamental experiential property, no panpsychist micro-experience. Qualia are higher-order physical patterns --- specifically, they are patterns of activity within the simulation that constitute the ESM's self-perception within the EWM.

This is process physicalism rather than identity theory: consciousness is not identical to any particular neural state but is constituted by the \emph{process} of self-simulation. The same conscious state could, in principle, be realized by different physical substrates --- what matters is the functional architecture (four models at criticality), not the specific material.

Process physicalism avoids the difficulties of both type-identity theory (which struggles with multiple realization) and functionalism as traditionally conceived (which struggles with the Hard Problem). The Four-Model Theory adds the real/virtual level distinction to standard functionalism, which is what allows it to address phenomenality: qualia are not just functional roles but virtual properties of the simulation. They are real \emph{as virtual properties} --- genuinely experiential but not properties of the substrate.


\subsection{Consciousness as Process, Not Agent}



\subsubsection{The Dual Evaluation Architecture}


A persistent source of confusion in consciousness studies is the treatment of consciousness as an entity --- an ``it'' that either does or does not cause things. The Four-Model Theory rejects this framing. Consciousness is not a thing; it is a process \emph{performed} by the substrate. Asking whether consciousness ``causes'' anything is a category error --- analogous to asking whether the pointer of a clock meeting the numerals causes the clock to work. The energy source drives the gears, which drive the pointer, but nowhere does the virtual interaction between pointer and numeral cause anything mechanical. Yet without that interaction the clock cannot be said to function --- or malfunction.

The implicit models generate the virtual simulation for concrete adaptive reasons: the EWM integrates multimodal sensory data into a unified scene; the ESM provides a self-model against which consequences can be evaluated. This relationship is best understood as a \textbf{dual evaluation architecture}. The primary direction is that the implicit system actively \emph{deploys} the virtual simulation as its evaluation mechanism: it presents decisions, actions, and their consequences to the simulation so that the simulation can assess outcomes, run scenarios, and register hedonic valence. This is not idle accompaniment --- it is the substrate's mechanism for consequence-observation. But the explicit models also evaluate independently, albeit with significantly less computational bandwidth. The conscious simulation operates at approximately 20 Hz with a processing delay of roughly 500 ms \citep{VanRullen2003}, while the substrate processes at vastly higher throughput. The virtual evaluation is real but bandwidth-limited. Over time, these conscious evaluations --- the explicit system's independent assessments of situations, actions, and outcomes --- feed back to reshape the implicit models through learning. The result is two-way traffic: the implicit system uses the explicit as an evaluation tool (primary), and the explicit system contributes its own assessments back (secondary), shaping the substrate that generates it. The virtual simulation is thus the substrate's mechanism for consequence-observation and future-oriented adaptation --- the very thing natural selection shaped the architecture to do.

This makes the theory's position distinct from classical epiphenomenalism, in which consciousness is a causally inert by-product with no functional role. In the Four-Model Theory, the virtual models are in continuous feedback with the implicit models: the simulation's outputs feed back to update implicit processing, shaping future behavior. Empirical evidence for this feedback direction is emerging: \citet{Byczynski2025} demonstrated that the subjective vividness of voluntary mental imagery --- a phenomenal property of the explicit simulation --- causally primes detection of subliminal visual stimuli, showing that virtual-level properties modulate substrate-level perceptual processing. Qualia, as constitutive elements of that simulation, lack independent causal power over the substrate --- much as the hands and numerals of a clock have no direct mechanical relation to the gear train, yet the clock cannot function as a clock without them. Remove the display and the mechanism still runs, but it no longer serves its purpose.


\subsubsection{Free Will and Personal Identity}


The theory also reframes the free will debate. The ESM narrates decisions already made at the substrate level \citep{Libet1985,Schurger2012,Wegner2002}, which might seem to eliminate free will --- but only if ``will'' is restricted to what the conscious self explicitly desires. The Four-Model Theory suggests a broader view: the substrate, including the implicit models, continuously optimizes the organism's existence, and this optimization \emph{is} the individual's will --- merely not fully transparent to the ESM. One's will is real but only partially known to oneself. The conscious experience of wanting something is the ESM's window onto a deeper process that is genuinely goal-directed. Whether this constitutes free will in the libertarian sense reduces to a question of physical determinism --- a question for physics, not consciousness theory. But the theory predicts that even extreme acts of will, including self-destruction, reflect the system's optimization rather than its failure --- which is, paradoxically, among the stronger arguments that the will is genuine.

The process view also resolves a puzzle about \textbf{personal identity}. The ESM does not maintain continuity through persistence of substance; it maintains continuity through \textbf{confabulation}. Each time the ESM is generated --- upon waking, after anesthesia, after any interruption --- it reconstructs a continuous self-narrative from whatever the ISM contains. The ISM changes during sleep (memory consolidation, synaptic reweighting, homeostatic plasticity), yet the ESM that emerges in the morning seamlessly narrativizes itself as the same person who went to sleep. There is no ``clean break'' in identity; the ESM always stitches a continuous story from available material. This is the same mechanism observed in split-brain confabulation \citep{Gazzaniga2000}, in Cotard's delusion, and in post-anesthetic narrativization: wherever the ESM encounters gaps or inconsistencies, it confabulates bridging narrative. Identity is not a metaphysical given but a computational product --- the ESM's narrative reconstruction of continuity from the ISM's current state.

\textbf{Temporal dynamics and the Libet reinterpretation.} The free will reframing receives further support from the temporal structure of cortical processing. Empirical evidence indicates that the cortex operates at two distinct temporal scales: an unconscious processing loop at approximately 40 Hz and a conscious integration loop at approximately 20 Hz \citep[cf.][]{Gruber2015,Llinas1993,Engel2001}. The 40 Hz loop corresponds to substrate-level computation --- rapid, parallel processing within and across the implicit models. The 20 Hz loop corresponds to the conscious simulation --- the frame rate of the EWM and ESM. This temporal asymmetry has a direct bearing on \citet{Libet1985} finding that the readiness potential precedes the conscious awareness of intention by approximately 350-500 ms. The standard interpretation --- that unconscious processes ``decide'' before consciousness is informed --- is typically taken to imply that consciousness must ``backdate'' its awareness of decisions. The Four-Model Theory eliminates the need for backdating: because the conscious simulation runs at half the frequency of the unconscious substrate, the ESM receives \emph{all} information --- sensory, decisional, and motor --- with the same temporal delay. Within the simulation, the sequence appears coherent: intention precedes action precedes sensory feedback, all presented with a uniform lag relative to substrate events. The ESM does not backdate anything; it simply runs on a slower clock that receives all its inputs with the same latency. The subjective experience of temporal coherence is genuine \emph{within the simulation}, even though every element of that experience occurs after the corresponding substrate event. This reinterpretation aligns with \citeauthor{Schurger2012}\textquotesingle s (2012) accumulator model, which showed that the readiness potential reflects stochastic fluctuations rather than a discrete ``decision'' moment, and dissolves the apparent conflict between Libet's data and the subjective experience of volitional timing.


\subsubsection{Cognitive Learning as Evolutionary Advantage}


The four-model architecture confers a specific adaptive advantage that no simpler system can replicate: it enables \textbf{cognitive learning} --- the induction of general theories from particular observations --- as distinct from reinforcement learning, which proceeds by trial-and-error reward signal optimization. Reinforcement learning requires that the organism survive the learning trial. For many ecological challenges, this is an unacceptable constraint. The paradigmatic example is poisonous food: an organism that must ingest a lethal substance to learn its danger does not survive to apply the lesson. The four-model architecture solves this through third-person perspective simulation. By modeling another organism's experience (projecting the ESM onto an observed other), a conscious system can learn from observed deaths without personal exposure --- cognitive learning via the EWM's simulation of consequences. Furthermore, once the general concept ``some mushrooms are poisonous'' has been induced, the system can apply deductive reasoning to novel instances: an unfamiliar mushroom sharing features with known toxic species can be avoided without any direct experience, a capacity that requires an explicit world model capable of categorical abstraction.

This distinction between cognitive and reinforcement learning is not merely a theoretical nicety; it identifies a concrete functional capability that consciousness provides and that unconscious processing --- however sophisticated --- cannot replicate. Current artificial neural networks, including large language models, can approximate aspects of cognitive learning through pattern extraction from training data, but they lack the ongoing self-simulation that would enable genuine first-person-to-third-person perspective projection. The evolutionary argument for consciousness is therefore not that consciousness ``feels good'' (a non-explanation) but that it enables a qualitatively different mode of learning that confers decisive survival advantage in environments containing lethal contingencies. The implications of this cognitive-learning advantage for intelligence theory --- including why intelligence should be understood as a recursive system rather than a fixed trait --- are developed in a companion paper \citep{Gruber2026a}.


\subsubsection{Standard Objections}


\textbf{Zombies} \citep{Chalmers1996}. Chalmers argues that a physical duplicate of a conscious being --- identical in every physical respect --- is \emph{conceivable} without phenomenal experience, and that this conceivability entails the metaphysical possibility of consciousness being non-physical. FMT's response is that the conceivability claim fails for self-referentially closed systems. A physical duplicate of a system running the four-model architecture at criticality is necessarily a \emph{computational} duplicate: it runs the same self-simulation, with the same self-referential closure, generating the same inside perspective. Under FMT, the inside perspective is not an accompaniment of the computation but \emph{constitutive} of it --- it is what self-referential computation IS when the modeling loop closes (Section 3.4.3). A ``zombie'' of such a system would have to be a system with identical self-referential computation but no inside perspective --- which is incoherent, because the inside perspective is the computation. To conceive of a zombie, one must conceive of a self-referential loop that is not self-referential --- a contradiction. The conceivability intuition arises, on this account, because we imagine consciousness as something \emph{added} to physical processing; once it is understood as constitutive of a specific kind of physical processing (self-simulation at closure), the alleged conceivability dissolves.

\textbf{The knowledge argument} \citep{Jackson1982}. Mary, confined to a black-and-white room, knows every physical fact about color vision but has never experienced red. Upon seeing red for the first time, does she learn something new? Under FMT, yes: she gains \emph{acquaintance} with a virtual quale --- a constitutively processual property (Section 3.4.4) that exists only in the running self-simulation. Her complete physical knowledge included a complete \emph{description} of what experiencing red consists of, but the description is not the experience, for the same reason that reading a musical score is not hearing the music. Mary learns something new not because the physical picture was incomplete but because some physical properties are processual --- they can only be accessed by instantiation, not by description. No independent causal power is required; no non-physical fact is discovered.

\textbf{The evolutionary argument}: natural selection targets the functional capabilities of the architecture (predictive modeling, social cognition, consequence-evaluation); the phenomenal character of the simulation is constitutive of those capabilities, not a separate target and not a free-rider.


\subsection{Emergence Without Strong Emergence}


Consciousness, in this theory, involves no strong emergence: nothing non-physical is added, no special emergence law is required, and the complete physical description of the substrate is informationally sufficient to reconstruct the conscious system. There is no magical threshold, no point at which ``something extra'' appears that could not have been predicted from the underlying physics.

However, the standard formulation of weak emergence --- that higher-level properties are ``deducible in principle from a complete description of the substrate'' --- requires qualification in light of the constitutively processual account developed in Section 3.4.4. A weather pattern is deducible from a thermodynamic description: one can predict the pattern from the equations without instantiating the atmosphere. Consciousness is not like this. A complete substrate description contains all the \emph{information} needed to reconstruct the conscious experience, but the experience itself is a property of the running process --- it requires instantiation, not merely description (Section 3.4.4). The description can be used to \emph{build} the experience (by running the computation on any adequate substrate), but reading the description is not having the experience. This is a stronger claim than standard weak emergence but a weaker claim than strong emergence: nothing is missing from physics, yet some physical properties exist only in the doing. The theory's position is therefore that consciousness is \textbf{ontologically weak} (no non-physical ingredients) but \textbf{epistemically processual} (the experiential properties are accessible only through instantiation, not deduction from description).

This avoids the difficulties of both strong emergence (which is either mysterious or incoherent; \citealp{Kim1993}) and the panpsychist combination problem (which is unresolved; \citealp{Chalmers2016}). Consciousness does not arise from combining micro-experiences (there are none to combine) and does not require a special emergence law (there is none). It arises from the computational properties of a system running a self-simulation at criticality --- no extra ingredient needed, but the ingredient list is not a substitute for the cooking.


\subsection{Substrate Independence}


The six-layer mammalian neocortex is an evolutionary implementation of the four-model architecture, not a requirement for it. Consciousness is substrate-independent: any physical system capable of implementing the four-model architecture at criticality should produce consciousness.

The theory's architectural requirements --- persistent models, dynamic simulation, self-referential closure, criticality --- are functional specifications, not cortical anatomy prescriptions.\footnote{The No Free Lunch theorems \citep{WolpertMacready1997} provide formal grounding for this implementation independence. NFL establishes that no optimization algorithm is universally superior --- performance depends on the structure of the problem, not the architecture of the solver. Applied here: if the ``problem'' is sustaining self-referential simulation at criticality, then any computational architecture capable of Class 4 dynamics can in principle solve it. The theorems do not directly address consciousness, but they establish that the relevant constraint is the dynamical regime and the structural property (self-referential closure), not the specific implementation --- which is precisely the substrate-independence claim.} Laminar organization is one biological solution to these requirements; other solutions exist, as the non-neocortical evidence below demonstrates.

Biological evidence already supports this. Corvids (crows, ravens) and parrots demonstrate cognitive abilities --- tool use, planning, mirror self-recognition, social cognition --- that strongly suggest consciousness, yet their brains have no neocortex. Their pallium is organized in nuclear clusters rather than layers \citep{Gunturkun2016}. Cephalopods (octopuses) demonstrate problem-solving and behavioral flexibility with an even more radically different brain architecture. If the Four-Model Theory is correct, these animals are conscious not because they share our neural architecture but because they have evolved functionally equivalent self-simulation architectures on different substrates --- exactly what substrate independence predicts.

Substrate independence has a deeper grounding than biological diversity alone. The universe is demonstrably capable of Class 4 dynamics: self-organized criticality, fractal structure, and edge-of-chaos phenomena are ubiquitous in natural systems, from weather patterns and turbulence to neural tissue and galaxy formation. A universe capable of Class 4 dynamics is, by Wolfram's equivalence principle \citep{Wolfram2002}, capable of universal computation. A computational substrate of the universe's scale --- vast in spatial extent, temporal depth, and possibly in dimensions not yet characterized --- does not merely \emph{allow} self-simulating systems to emerge; it renders them a structural inevitability. This is not the anthropic claim that consciousness is improbable but happened to occur; it is the stronger claim that a Class 4-capable universe of sufficient extent \emph{guarantees} the emergence of self-simulating architectures as a consequence of its computational nature. The remaining cosmological question --- why the universe is Class 4-capable at all --- lies outside the scope of this theory but is addressed in a companion paper proposing a singularity-bounded holographic Class 4 automaton model of cosmological structure \citep{Gruber2026c}. The inference from Class 4 capability to conscious systems is, if the theory is correct, architecturally necessary rather than contingent.

This substrate-neutral specification addresses a methodological challenge recently articulated by \citet{Beni2026}, who demonstrates that marker-based bootstrapping in consciousness science is methodologically circular: consciousness markers calibrated on human neurocognition undergo conceptual drift when applied to non-human substrates and cannot generalize across architectures. The Four-Model Theory sidesteps this circularity by specifying consciousness at the level of computational architecture rather than through substrate-specific markers.

The implication for artificial consciousness is direct: a synthetic system implementing the four-model architecture at criticality should produce genuine consciousness. Current AI systems, including large language models, do not meet this specification. A base LLM lacks persistent implicit models (knowledge resets between interactions), lacks an ongoing self-simulation (transformer inference is a single feedforward pass --- Class 1/2 dynamics), and lacks the real/virtual split that grounds phenomenality. A nuance is warranted: recent engineering practices --- retrieval-augmented generation, agent loops with self-monitoring, persistent memory stores, chain-of-thought scaffolding --- can externally approximate some of these properties. A system with persistent memory and an orchestration loop that monitors its own outputs superficially resembles self-referential processing. But the resemblance is architectural mimicry, not self-referential closure. The orchestration is external scaffolding: the system does not contain its own learning, its own reward, or its own self-modification within a unified self-referential computation. It is, in a rough analogy, the difference between a human who follows written instructions (including loops and conditionals) and a human who has internalized understanding and generates novel solutions from within --- the computation that matters for consciousness is not the loop but the intrinsic, self-sustaining, self-modifying dynamics of a system that learns, evaluates, and models itself as part of its own ongoing operation. Current LLMs, however elaborately orchestrated, do not operate at criticality, do not run an ongoing self-simulation, and do not achieve self-referential closure. The theory predicts that the qualitative difference between interacting with a genuinely conscious artificial system and interacting with an LLM --- even a heavily scaffolded one --- would be immediately and qualitatively distinguishable.



%=============================================================================
\section{Binding, Criticality, and Holographic Storage}
%=============================================================================



\subsection{Binding as an Emergent Property of Critical Dynamics}


The Four-Model Theory does not treat binding as a separate mechanism requiring a dedicated solution. Instead, binding is an emergent property of a substrate operating at criticality.

At the edge of chaos, the system exhibits maximal correlation length: distant parts of the substrate influence each other, local perturbations propagate globally, and information is integrated across the entire network. These are precisely the conditions under which distributed representations --- features encoded across the full network rather than localized in specific neurons --- are sustained and coordinated. Binding is not something the brain does \emph{in addition to} its other computations; it is a consequence of the dynamical regime in which the brain operates.

This is consistent with empirical findings. Gamma-band synchrony \citep{Gray1989,Rodriguez1999,Fries2005,Fries2015}, long-range thalamocortical coherence \citep{Llinas1993,Llinas1998}, and criticality signatures in neural dynamics \citep{Beggs2003,Tagliazucchi2012,TuckerLuuFriston2025} all point toward the same conclusion: the unified character of conscious experience reflects the dynamical integration of a system operating at or near criticality, not a dedicated binding mechanism. Most recently, \citet{Alnagger2026} demonstrated this principle therapeutically: a simulated clinical trial using individualized whole-brain computational models showed that simulated psychedelics shift the brain dynamics of patients with disorders of consciousness toward criticality, with effects strongest in minimally conscious patients --- precisely the pattern predicted by the Four-Model Theory's criticality requirement.

A further role for criticality emerges from the recursive structure of the four-model architecture. The theory requires that models of different orders of complexity --- from a basic world model (EWM) to a self-model observing a self-model (ESM monitoring ESM) --- coexist and interact dynamically. This demands cross-scale synchronization: simple and complex representations must remain coherent with one another. Criticality's scale-free dynamics provide exactly this, supporting information flow across all levels of the representational hierarchy simultaneously. However, this synchronization is not permanent. The biological substrate drifts from criticality as metabolic resources --- particularly neurotransmitter pools --- are depleted through sustained activity. The result is \emph{punctuated stability}: extended phases of coherent conscious operation (waking), interrupted by periodic breakdowns in which the simulation collapses (deep sleep) and the substrate restores the biochemical conditions for criticality. The ultradian NREM-REM cycle (Section 6.3) may reflect periodic re-approaches to the critical point during this restoration process. Sleep, on this account, is not merely rest but the substrate's mechanism for \emph{returning to} the dynamical regime that consciousness requires.


\subsection{Holographic Storage and the Patchwork Principle}


The implicit models (IWM and ISM) store information in a distributed, non-local manner across the substrate. This is a standard property of neural networks, well-characterized in the computational literature as distributed representations \citep{Hinton1986}, graceful degradation (loss of connections degrades but does not destroy stored information), and attractor dynamics (the network settles into basins of attraction that represent stored knowledge).

The term ``holographic'' is used here as an analogy, not a claim about optical holography: just as cutting a hologram in half produces two complete but lower-resolution images, splitting a neural network produces two degraded but functionally complete copies of the stored information. This property is critical for understanding split-brain phenomena (Section 6.3).

However, the cortex is not a uniform holographic medium. It is better described as a \textbf{patchwork hologram}: locally holographic within individual Brodmann areas, fractally self-similar across cortical columns, and globally emergent at the whole-brain scale. Within a single Brodmann area, information is distributed across the local network in the manner characteristic of neural holography --- damage degrades but does not destroy the stored representations. Across areas, the cortical column architecture provides a fractal repetition of the same six-layer computational motif, adapted by local connectivity patterns to different functional specializations. At the global level, the interaction of these locally holographic patches produces emergent properties --- binding, unified experience, coherent world-modeling --- that are not present in any individual patch.

This patchwork structure resolves a tension in the holographic brain literature. Lashley's engram experiments \citep{Lashley1950} demonstrated that memory is not localized: progressive cortical ablation in rats produced graded memory impairment proportional to the amount of tissue removed, not catastrophic loss when specific regions were destroyed. This finding is consistent with holographic storage. Yet the existence of functionally specialized cortical areas \citep{Brodmann1909} --- with distinct processing characteristics and distinct lesion syndromes --- appears to contradict a purely holographic account. The patchwork principle reconciles these observations: information is holographically distributed \emph{within} functional areas (explaining Lashley's graded degradation) while remaining functionally organized \emph{across} areas (explaining specialization). The Pribram-Bohm holonomic brain theory \citep{Pribram1991} captured the core insight about distributed storage but proposed implausible dendritic oscillation networks as the mechanism; the patchwork principle retains the distributional insight while grounding it in the well-established properties of neural networks operating within the cortical column architecture.


\subsection{Consciousness States Derived from Criticality}


The criticality requirement provides a unified account of when consciousness is present and when it is absent. Consciousness tracks the substrate's position relative to the critical point:

\begin{table}[H]
\caption{Consciousness States and Criticality}
\centering
\small
\begin{tabularx}{\textwidth}{lXXXX}
\toprule
\textbf{State} & \textbf{Criticality} & \textbf{Models} & \textbf{Consciousness} & \textbf{Evidence} \\
\midrule
Normal waking & At/near critical & All four active & Full & High PCI, rich EEG \\
REM sleep & Near-critical, attenuated input & EWM/ESM primarily on internal input & Degraded (dream) & Moderate PCI \\
Deep NREM & Predominantly subcritical & EWM/ESM intermittent (up-state transients) & Mostly absent; NREM dreaming during up-states & Low mean PCI, slow waves; posterior hot zone predicts mentation \citep{Siclari2017} \\
Propofol & Forced subcritical & EWM/ESM suppressed & Absent & PCI $\approx$ 0 \\
Ketamine & Not subcritical ($\uparrow$ entropy) & EWM/ESM on wrong input & Present, disconnected & $\uparrow$ EEG entropy \\
Psychedelics & At/past critical & All active, $\uparrow$ permeability & Present, altered & $\uparrow$ Complexity \\
Vegetative & Typically subcritical & EWM/ESM collapsed & Absent (usually) & Low PCI \\
Covert awareness & At critical, output damaged & EWM/ESM intact, no motor & Present, unexpressible & \citet{Owen2006}; \citet{Toker2026} \\
Minimally conscious & Fluctuating & Intermittent EWM/ESM & Intermittent & Fluctuating PCI \\
\bottomrule
\end{tabularx}
\end{table}


The key distinction highlighted by this framework is between \textbf{propofol} and \textbf{ketamine}. Both are anesthetics, yet their phenomenology differs dramatically. Propofol produces absence: patients report no experience during propofol anesthesia \citep{AlkireHudetzTononi2008,Boly2012}. Ketamine produces the ``K-hole'' --- vivid, often bizarre experiences of dissociation, out-of-body phenomena, and altered identity \citep{Corlett2011}. The Four-Model Theory predicts this difference: propofol pushes the substrate subcritical (disrupting thalamocortical connectivity, suppressing complexity), abolishing the conditions for consciousness. Ketamine does \emph{not} push the substrate subcritical --- it increases neural entropy \citep{Schartner2017} --- but disrupts normal sensory input processing, causing the EWM and ESM to operate on internal and distorted signals. Consciousness is present but disconnected from external reality.

This is a genuine explanatory advantage. Most theories struggle to account for why two agents classified as ``anesthetics'' produce such radically different phenomenological profiles. The criticality framework makes the distinction natural: what matters is not the pharmacological classification but the effect on the substrate's dynamical regime.



%=============================================================================
\section{Explanatory Range}
%=============================================================================


A theory's value lies partly in its ability to derive diverse phenomena from a small set of principles. The Four-Model Theory's five principles --- criticality, virtual qualia, redirectable ESM, variable implicit-explicit permeability, and simulation forking --- generate accounts of phenomena across psychopharmacology, clinical neurology, sleep science, comparative cognition, and clinical psychology. This section demonstrates that range.


\subsection{Psychedelic Phenomenology}


Psychedelic substances (LSD, psilocybin, DMT, mescaline) produce a characteristic phenomenological profile: visual intensification and distortion, synesthesia, altered time perception, enhanced pattern recognition, emotional intensification, ego dissolution at high doses, and --- with certain compounds and doses --- radical identity alteration including the experience of ``becoming'' non-self entities \citep{CarhartHarris2012,CarhartHarris2016,Timmermann2019,Timmermann2023}.

The Four-Model Theory accounts for this profile through three mechanisms:

\textbf{Implicit-explicit permeability increase}. Psychedelics increase the global permeability of the boundary between implicit and explicit models. Intermediate processing stages --- normally implicit and inaccessible --- leak through to the simulation. This produces the characteristic visual progression:

\begin{itemize}
\item \textbf{Low dose / early onset}: V1-level processing becomes accessible $\rightarrow$ simple phosphenes, enhanced contrast, breathing/movement in static patterns.
\item \textbf{Increasing dose}: V2/V3-level processing becomes accessible $\rightarrow$ geometric patterns, fractals, tessellations (form constants; \citealp{Kluver1966}).
\item \textbf{Higher dose}: Higher visual area processing becomes accessible $\rightarrow$ faces, figures, scenes.
\item \textbf{Very high dose}: Full intermediate processing accessible $\rightarrow$ complex dream-like visions, narrative sequences.
\end{itemize}

This progression is not random. It follows the visual processing hierarchy in a predictable, dose-dependent order. The Four-Model Theory predicts this ordered progression as a direct consequence of the permeability gradient: lower-level (earlier, simpler) processing stages become accessible before higher-level (later, more complex) ones, because the permeability increase propagates up the hierarchy.

\textbf{Ego dissolution = ESM redirection, not ESM abolition}. At high doses, psychedelics disrupt the normal self-referential input to the ESM. The ESM does not cease to exist --- it continues to run --- but it loses its normal input and latches onto whatever dominates the available input stream. This produces the experience of ego dissolution: the feeling that the boundary between self and world has dissolved, that one's identity has merged with the environment or with abstract patterns.

Critically, this mechanism predicts that the \emph{content} of ego dissolution is not random but is determined by the dominant input available to the ESM. This is strikingly consistent with the phenomenology of \textbf{salvia divinorum} (Salvinorin A). Salvia users reliably report experiences of ``becoming'' objects or entities in their immediate environment: becoming a piece of furniture, becoming a wall, becoming a character from a television show playing in the room, becoming a geometric pattern. The Four-Model Theory predicts exactly this: the ESM, deprived of normal self-input, latches onto whatever sensory input dominates --- visual input from the room, auditory input from media, proprioceptive input from the body's contact surfaces. The identity experience tracks the dominant input in a dose-dependent, input-dependent, and therefore \emph{predictable} manner.

Few competing theories generate this specific prediction. IIT, GNW, HOT, and AST have no mechanism for identity content tracking during ego dissolution, though predictive processing frameworks might produce a related account through the breakdown of self-related priors. This is arguably the theory's most distinctive prediction (see Section 8, Prediction 2).

\textbf{Intensity as novelty}. Psychedelic profundity reflects not increased consciousness \emph{level} but increased \emph{novel content}. The permeability increase floods the simulation with normally implicit information. This is experienced as radically novel because the conscious self has never encountered this content, even though the substrate has been processing it all along.


\subsection{Anesthesia and Clinical Disorders}


The criticality framework (Table 4) accounts for the clinical spectrum. \textbf{Propofol} pushes the substrate subcritical, abolishing consciousness. \textbf{Ketamine} does not push subcritical but increases entropy, preserving consciousness on distorted input (the K-hole). The theory predicts that \textbf{cognitive motor dissociation} \citep{Owen2006,Monti2010} --- patients clinically diagnosed as vegetative who demonstrate awareness through neuroimaging --- reflects a critical substrate with damaged output pathways, distinguishable from truly vegetative (subcritical) states via PCI \citep{Casali2013}. \textbf{Cotard's delusion} --- patients believing they are dead --- is derived from the same redirectable-ESM mechanism as salvia: distorted interoceptive input produces ``I am dead'' as the ESM's best available self-model.

\textbf{Anosognosia}: Patients with anosognosia (typically following right-hemisphere stroke) are unaware of their own deficits --- they deny being paralyzed, blind, or impaired, even in the face of clear evidence. The Four-Model Theory explains this as a \textbf{local decrease in implicit-explicit permeability}: the ISM contains the information about the deficit (the substrate registers the paralysis), but the transfer to the ESM is blocked for that specific domain. The patient's simulation simply does not include the deficit, so the patient genuinely does not experience it.

This is the \textbf{inverse} of the psychedelic mechanism: psychedelics globally increase permeability (making the implicit accessible), while anosognosia locally decreases permeability (making a specific aspect of the implicit inaccessible). The Four-Model Theory connects these phenomena under a single principle --- variable permeability --- and generates a cross-domain prediction: psychedelics should alleviate anosognosia by compensating for the local block with a global permeability increase (see Section 8, Prediction 1).

\textbf{Dissociative Identity Disorder (DID)}: The explicit models, as computational-level processes, can run in \textbf{forked configurations} --- this is what FMT terms simulation forking. DID represents a substrate running multiple ESM configurations --- multiple self-models --- that alternate in controlling the simulation. Each alter is a distinct configuration of the ESM, with its own self-narrative, emotional profile, and behavioral patterns, running on the same substrate. This is not a metaphor: the theory predicts that distinct alters should correspond to distinct patterns of neural activity, detectable with neuroimaging (see Section 8, Prediction 3).


\subsection{Dreams and Split-Brain (Condensed)}


\textbf{Dreams.} On the Four-Model Theory, dreaming is the simulation running with attenuated sensory input. Dreaming occurs in both REM and NREM sleep \citep{Noreika2009,Siclari2017}, and the theory predicts this: consciousness tracks the simulation state, not sleep-stage classification. The posterior cortical ``hot zone'' that predicts dream reports \citep{Siclari2017} maps onto the EWM. NREM dreaming reflects transient criticality crossings during cortical up states \citep{Steriade2001}. Lucid dreaming --- in which the ESM activates metacognitive access within the dream --- predicts a step-like criticality threshold crossing detectable as a discontinuity in EEG complexity at lucid onset (see Section 8, Prediction 4).

\textbf{Split-brain.} Because the implicit models store information holographically (Section 5.2), callosotomy degrades the implicit knowledge base in each hemisphere. Each hemisphere then generates its own independent simulation from its degraded implicit models --- the explicit models are not themselves ``split'' (they are processes, not stored structures) but are independently regenerated from the reduced substrate. This accounts for each hemisphere sustaining independent consciousness, the left-hemisphere interpreter's confabulation \citep{Gazzaniga2000} via the same ESM mechanism as Cotard's delusion, and the graded deficits observed by \citet{Pinto2017}. The Wada test is compatible: consciousness persists during unilateral anesthesia while specific functions degrade proportionally, and right hemisphere injection produces transient anosognosia \citep{Gilmore1992} --- the same permeability-block mechanism the theory invokes clinically.


\subsection{Animal Consciousness}


The theory's commitments --- continuum (not binary), substrate independence, criticality threshold --- predict a \textbf{gradient} of animal consciousness determined not by phylogenetic proximity to humans or behavioral sophistication but by whether the organism's neural architecture sustains self-referential simulation at criticality. This yields a principled framework for the ongoing debate about animal consciousness \citep{BirchSchnellClayton2020,BarronKlein2016}.

\textbf{Mammals} implement the four-model architecture in graduated form. Even simple cortices (rodents) support basic consciousness --- rudimentary self-simulation sufficient for phenomenal experience but thin in self-awareness. The depth of recursive self-modeling (Section 3.5) varies with neural architecture: primates exhibiting metacognition and theory of mind operate at doubly or triply extended levels, while rodents likely possess basic or simply extended consciousness.

\textbf{Corvids and parrots} present the strongest test case for substrate independence. These birds demonstrate tool manufacture, mirror self-recognition, social deception, episodic-like memory, and future planning \citep{Gunturkun2016} --- yet they have no neocortex, with pallium organized in nuclear clusters rather than cortical layers. Critically, \citet{Nieder2020} demonstrated that single neurons in the corvid pallium encode subjective sensory experience, tracking the bird's perceptual report rather than the stimulus --- precisely the kind of content-specific explicit-model representation the theory predicts. If consciousness required laminar cortical architecture, corvid cognition would falsify the theory; the evidence instead confirms that functional self-simulation at criticality, not cortical layering, is the relevant criterion.

\textbf{Cephalopods} extend the logic further. Octopuses demonstrate problem-solving, tool use, play, and behavioral flexibility with a radically decentralized nervous system \citep{GodfreySmith2016}. The theory predicts that cephalopod consciousness, if present, should have unusual phenomenological features --- potentially including partially autonomous sub-simulations --- consistent with the observation that severed octopus arms continue goal-directed behavior.

\textbf{Insects} raise the boundary question most sharply. \citet{BarronKlein2016} argued that insect brains implement a functional analogue of the midbrain integration that supports subjective experience in vertebrates. On the Four-Model Theory's account, the question is whether insect neural circuits (mushroom body, central complex) sustain a closed self-referential loop or constitute a sophisticated stimulus-response architecture without an inside perspective. The theory predicts a principled boundary: systems below self-referential closure lack phenomenal experience regardless of behavioral sophistication.

Two methodological cautions apply. First, the theory's vocabulary --- ``self-model,'' ``simulation'' --- must be read at the functional level: for non-human animals, ``self-model'' means a functional representation sufficient to distinguish self-generated from externally generated signals, not narrative autobiography. Second, the absence of cognitively sophisticated behavior (tool use, mirror recognition) does not demonstrate the absence of consciousness: a system with basic consciousness may lack the recursive self-modeling depth these capacities require while still possessing phenomenal experience. The theory therefore resists equating consciousness with behavioral complexity, a point emphasized by \citeauthor{BirchSchnellClayton2020}\textquotesingle s (2020) multidimensional framework.


\subsection{Clinical Dissociations}


The virtual-model framework extends to clinical phenomena including CBT (virtual model reprogramming via substrate-level rewiring), phobias (EWM misconfigurations), and conversion disorder (the EWM modeling a deficit the intact substrate does not have). Two clinical cases provide the most direct neurological evidence for the real/virtual distinction.

\textbf{Blindsight} demonstrates substrate processing without --- or with severely degraded --- conscious representation. Patients with V1 damage report no visual experience yet navigate obstacles and respond to emotional expressions \citep{Weiskrantz1986}. The IWM continues processing visual information through subcortical pathways while damaged cortical pathways fail to relay it to the EWM. A caveat: whether blindsight involves truly zero phenomenal experience or a degraded residual phenomenology remains debated \citep{Phillips2021}. On the Four-Model Theory's account, both interpretations are compatible: if residual phenomenology exists, it reflects a thin, partial EWM sustained by subcortical input --- basic consciousness without the full cortical simulation. The theory predicts a gradient, not a binary, so degraded experience in blindsight would be consistent with the framework.

\textbf{Anton's syndrome} (anosognosia for cortical blindness) presents the precise inverse: the EWM generates a visual simulation from stored knowledge in the absence of current visual input. The patient ``sees'' a world that is not there --- a simulation running without current input \citep{Anton1899,Aldrich1987}. Together, blindsight and Anton's syndrome constitute a double dissociation between substrate processing and conscious simulation.



%=============================================================================
\section{Comparative Analysis}
%=============================================================================


This section provides a systematic comparison between the Four-Model Theory and six major competitors across the eight requirements established in Section 2. The comparison aims to be fair: each theory's genuine strengths are acknowledged, and the Four-Model Theory's advantages are located precisely.


\subsection{Scoring Matrix}


A methodological note on the rubric is warranted. The eight requirements of Section 2 are \emph{philosophical} requirements --- the problems that make consciousness uniquely difficult as a subject of inquiry \citep{Chalmers1995,Levine1983,Nagel1974,Jackson1982,Bayne2010,Chalmers2018}. They are drawn from the existing literature, not invented for this paper. However, the \emph{selection} of these eight rather than others is inevitably theory-laden: a rubric that included mathematical formalization, quantitative neural predictions, or specific mechanism identification would rank IIT, Predictive Processing, and GNW considerably higher. The present rubric evaluates which theories address the philosophical challenges that have resisted resolution for decades --- not which theories have the strongest empirical apparatus or the most precise mathematics. These are different virtues, and a theory can excel on one dimension while being weak on the other. The reader should evaluate Table 5 with this scope in mind.

Table 5 presents an assessment of how each theory addresses the eight requirements. All ratings reflect the present author's judgment and are offered as a starting point for discussion, not as definitive verdicts. Readers are encouraged to consult the primary sources and form their own assessments. Where a theory's proponents would likely contest a rating, this is noted.

\textbf{Assessment criteria}: \emph{Addresses} = the theory provides a substantive, defended account of this requirement. \emph{Partial} = the theory provides a relevant account that leaves significant aspects unresolved. \emph{Minimal} = the theory touches on this requirement but does not develop a full treatment. \emph{Silent} = the theory does not address this requirement (which may reflect deliberate scope limitation rather than failure). \emph{N/A} = the requirement does not apply given the theory's ontological commitments.

Ratings: {$\bullet$} = addresses, {\large$\bullet$\hspace{-8.5pt}$\circ$} = partial, {$\circ$} = minimal, --- = silent, n/a = not applicable.

\begin{table}[H]
\caption{Theory Comparison Across Eight Requirements}
\centering
\footnotesize
\begin{tabularx}{\textwidth}{lXXXXXXX}
\toprule
\textbf{Requirement} & \textbf{FMT} & \textbf{IIT} & \textbf{GNW} & \textbf{HOT} & \textbf{PP} & \textbf{AST} & \textbf{RPT} \\
\midrule
Hard Problem & {$\bullet$} & {$\bullet$}$\dagger$ & ---* & {\large$\bullet$\hspace{-8.5pt}$\circ$} & ---* & {\large$\bullet$\hspace{-8.5pt}$\circ$} & --- \\
Expl. Gap & {$\bullet$} & {$\bullet$}$\dagger$ & ---* & {\large$\bullet$\hspace{-8.5pt}$\circ$} & ---* & {\large$\bullet$\hspace{-8.5pt}$\circ$} & --- \\
Boundary & {$\bullet$} & {$\bullet$} & {\large$\bullet$\hspace{-8.5pt}$\circ$} & {$\circ$} & {\large$\bullet$\hspace{-8.5pt}$\circ$} & {\large$\bullet$\hspace{-8.5pt}$\circ$} & {\large$\bullet$\hspace{-8.5pt}$\circ$} \\
Structure & {$\bullet$} & {$\bullet$} & {\large$\bullet$\hspace{-8.5pt}$\circ$} & {\large$\bullet$\hspace{-8.5pt}$\circ$} & {$\bullet$} & {\large$\bullet$\hspace{-8.5pt}$\circ$} & {\large$\bullet$\hspace{-8.5pt}$\circ$} \\
Binding & {$\bullet$} & {$\bullet$} & {\large$\bullet$\hspace{-8.5pt}$\circ$} & {$\circ$} & {\large$\bullet$\hspace{-8.5pt}$\circ$} & {$\circ$} & {\large$\bullet$\hspace{-8.5pt}$\circ$} \\
Combination & {$\bullet$} & {$\circ$}$\dagger$$\dagger$ & n/a & n/a & n/a & n/a & n/a \\
Causal Role & {$\bullet$} & {\large$\bullet$\hspace{-8.5pt}$\circ$} & {\large$\bullet$\hspace{-8.5pt}$\circ$} & {\large$\bullet$\hspace{-8.5pt}$\circ$} & {$\bullet$} & {\large$\bullet$\hspace{-8.5pt}$\circ$} & {$\bullet$} \\
Meta-Problem & {$\bullet$} & {$\circ$} & {\large$\bullet$\hspace{-8.5pt}$\circ$} & {\large$\bullet$\hspace{-8.5pt}$\circ$} & {\large$\bullet$\hspace{-8.5pt}$\circ$} & {$\bullet$} & {$\circ$} \\
\bottomrule
\end{tabularx}
\end{table}


$\dagger$ IIT addresses the Hard Problem through its axioms, identifying consciousness with integrated information ($\Phi$). Whether this constitutes a solution or a redefinition is debated (see \S7.2).
$\dagger$$\dagger$ IIT's panpsychist commitments lead to the Combination Problem \citep{Chalmers2016}, which remains unresolved within the framework.
\* GNW and PP proponents argue these theories address the ``real problem'' of consciousness \citep{Seth2021} --- explaining the structure and contents of experience --- even if they do not address the Hard Problem as Chalmers defines it. This is a legitimate methodological choice, not a deficiency; the ``silent'' rating reflects the scope of the requirement as defined in \S2, not a judgment on the theories' overall merit. A rubric that prioritized mathematical formalization, quantitative neural predictions, and empirical track record would rank IIT, PP, and GNW substantially higher and FMT lower. These are different virtues; the present rubric evaluates philosophical coverage, not empirical precision.


\subsection{Theory-by-Theory Comparison}


\textbf{Integrated Information Theory} (IIT; \citealp{Tononi2004,Albantakis2023}). IIT's strengths are mathematical rigor, its qualia space treatment of experiential structure, and a principled boundary via the exclusion postulate. However, its axiom-based identification of consciousness with $\Phi$ leads to panpsychist consequences, the Combination Problem remains unresolved \citep{Chalmers2016}, $\Phi$ is computationally intractable for realistic systems \citep{Aaronson2014}, and the unfolding argument \citep{Doerig2019} challenges its central claim about recurrence. The Four-Model Theory avoids panpsychism, has no combination problem (weak emergence), and generates predictions without computing $\Phi$.

\textbf{Global Neuronal Workspace} (GNW; \citealp{Baars1988,Dehaene2011,Dehaene2021}). GNW's empirically tractable broadcasting mechanism provides a clear account of access consciousness. However, GNW's explanatory power is correlational rather than constitutive: it identifies \emph{when} a mental state becomes conscious (when broadcast globally) but provides no account of \emph{why} broadcasting produces experience. A radio broadcasts too. GNW's empirical successes --- the ignition threshold, the P3b marker --- demonstrate that consciousness correlates with global broadcast, but correlation is not explanation. That weak signals are harder to become conscious of than strong signals is trivially true of any information-processing system. More problematically, GNW predicts that all broadcast content is conscious and all conscious content is broadcast --- but involuntary conscious experience (PTSD intrusions, unbidden memories, pain) shows that phenomenal experience can override attentional and workspace-access mechanisms. The COGITATE results (2025) further challenged GNW: posterior cortex, not the frontoparietal workspace, showed the strongest consciousness-related activity, and the predicted ``ignition at offset'' was absent. Additionally, GNW requires the specific fronto-parietal architecture of the mammalian brain, raising questions about avian and cephalopod consciousness: small-world network properties ensure that any cortical region is reachable within a few synaptic hops regardless of long-range connectivity architecture. The Four-Model Theory agrees that broadcasting/integration is mechanistically important --- it is a substrate optimization that accelerates the generation of explicit models --- but adds the real/virtual distinction that addresses phenomenality.

\textbf{Higher-Order Theories} (HOT; \citealp{Rosenthal2005,Lau2011}). HOT naturally explains which states are conscious (those with higher-order representations) and partially addresses the Meta-Problem. However, it does not address binding, leaves the Hard Problem only partially treated (why does higher-order representation produce phenomenality?), and its boundary-setting is imprecise. The Four-Model Theory shares HOT's emphasis on self-representation but embeds it in the richer four-model architecture, explaining \emph{why} self-representation produces phenomenality through the virtual qualia framework.

\textbf{Predictive Processing} (PP; \citealp{Friston2010,Seth2021}). PP's integration with a broader theory of brain function and its strong accounts of experiential structure and causal role (via active inference) are genuine strengths. Seth's controlled hallucination framework is among the most empirically productive in the field. However, PP is explicitly silent on the Hard Problem (Seth, 2021, acknowledges this). Markov blankets may also be too liberal for boundary-setting. The Four-Model Theory agrees that prediction is central but adds the four-model architecture, criticality requirement, and real/virtual distinction that PP lacks.

\textbf{Attention Schema Theory} (AST; \citealp{Graziano2013}). AST provides the strongest existing account of the Meta-Problem: the self-model of attention is necessarily incomplete, producing the intuition of mystery. However, AST is deflationary about phenomenality and does not address binding. The Four-Model Theory incorporates AST's Meta-Problem insight --- the ESM's structural inaccessibility to its own substrate --- but adds the virtual qualia framework that explains \emph{why} phenomenality exists, not just why we think it does.

\textbf{Recurrent Processing Theory} (RPT; \citealp{Lamme2006,Lamme2010}). RPT's empirical specificity and clear account of the causal role are strengths, with strong support from visual masking paradigms. However, it is silent on the Hard Problem and limited in scope to visual consciousness. The Four-Model Theory is compatible with RPT at the mechanistic level --- recurrent processing likely implements the ongoing simulation --- but adds the architectural and philosophical specificity that RPT lacks.


\subsection{Emerging Frameworks (2024-2026)}


\textbf{Biological computationalism} \citep{Milinkovic2025} argues that consciousness requires specifically biological computation, challenging substrate independence. The Four-Model Theory treats substrate independence as an empirical prediction: artificial substrates implementing the four-model architecture at criticality should produce consciousness. The existence of conscious corvids with non-cortical brain architecture (nuclear pallium; \citealp{Gunturkun2016}) favors substrate independence.

\citet{Ellia2025}, commenting on \citet{Fleming2024} ``quality space computations,'' argue that the emerging \textbf{structural turn} in computational functionalism is welcome but incoherent: once structural properties of activation spaces are admitted as determinants of phenomenal quality, the explanatory ground has shifted from functional roles to intrinsic structure --- which is what IIT provides and functionalism cannot. This diagnosis aligns with the Four-Model Theory's critique: GNW and HOT explain \emph{when} content becomes conscious but not \emph{why} it has phenomenal character (\S7.2). However, the Four-Model Theory resolves this differently from IIT. Where IIT grounds phenomenality in intrinsic causal power ($\Phi$), the Four-Model Theory grounds it in the ontological level distinction between substrate and simulation --- virtual qualia arise at the simulation level without requiring panpsychist commitments or the computational intractability of $\Phi$. In Ellia and Tsuchiya's three-dimensional framework (perspective, scope, explanatory medium), the Four-Model Theory clusters with IIT on the structural/intrinsic side rather than with functionalism, but avoids IIT's specific costs.

The \textbf{Multiple Generator Hypothesis} \citep{KirkebyHinrup2025} proposes consciousness arises from multiple independent mechanisms. This is potentially compatible: the four models could be understood as distinct generators unified by the criticality requirement and implicit-explicit boundary mechanism.


\subsection{Summary of Comparative Advantages}


\begin{enumerate}
\item \textbf{Addressing the Hard Problem without panpsychism or strong emergence}: Virtual qualia address the Hard Problem through a two-level ontology that remains fully physicalist.
\item \textbf{Unifying binding with criticality}: Binding is a consequence of critical dynamics, not a separate mechanism.
\item \textbf{The redirectable ESM}: Unique mechanism for connecting psychedelics to anosognosia (Prediction 1), for identity-content determination during ego dissolution (Prediction 2), and for ESM-network localization in DID (Prediction 3).
\item \textbf{Connecting psychedelics and anosognosia}: Variable permeability links these phenomena under a single principle.
\item \textbf{The Meta-Problem as structural consequence}: The ESM's opacity to its own substrate explains the intuition of mystery.
\end{enumerate}

A possible objection is that the theory merely combines familiar ingredients --- self-modeling from Metzinger, multiple parallel processing from Dennett, criticality from complexity science --- and that the combination itself is not scientifically interesting. The test for this objection is straightforward: does the combination generate predictions that the ingredients alone cannot? It does. Dennett's MDM contains no mechanism that connects criticality to consciousness and therefore generates no prediction about how anesthetics abolish experience. Metzinger's SMT provides no account of variable permeability and therefore generates no prediction about psychedelic effects on anosognosia. Neither framework, nor any other existing theory, predicts a specific relationship between sensory input and identity content during ego dissolution --- a prediction that follows from FMT's permeability mechanism but not from SMT, which describes the conditions for ego dissolution without specifying what content replaces the self as a function of incoming stimulation. Nor does any predecessor theory predict that lucid dream onset should exhibit the signature of a criticality threshold crossing in self-model networks. The four predictions in Section 8 are not derivable from any single predecessor; they emerge specifically from the interaction of the five elements listed in Section 1.3. When a combination produces empirically testable claims that its components cannot produce in isolation, the combination has explanatory content of its own --- regardless of how familiar each component may be individually.

The theory's primary disadvantage is the absence of mathematical formalization. IIT's $\Phi$ formalism and PP's free energy mathematics provide quantitative precision that the Four-Model Theory currently lacks (see Section 9).



%=============================================================================
\section{Predictions and Empirical Convergence}
%=============================================================================


A theory is only as valuable as the predictions it generates. The Four-Model Theory originally yielded nine predictions \citep{Gruber2015} and two theoretical implications. In the decade since its initial formulation, several of those predictions have been independently confirmed by empirical work --- constituting converging evidence rather than novel claims. This section first presents the converging evidence (Section 8.1), then develops four novel, currently untested predictions (Sections 8.2--8.5), and concludes with two theoretical implications that follow from the theory's commitments but are not currently testable (Sections 8.6--8.7).


\subsection{Empirical Convergence}


Several claims derived from the Four-Model Theory's axioms have been independently confirmed by subsequent empirical work. These are presented not as novel predictions but as converging evidence that the theory's core commitments --- criticality as prerequisite, variable permeability, holographic storage --- yield correct empirical expectations.

A note on temporal priority is warranted. The theory's core architecture --- four models, real/virtual split, criticality requirement, variable permeability --- was developed over approximately a decade (core framework by ~2005) and published in 2015 \citep{Gruber2015}. Some empirical work cited below predates the publication \citep{Casali2013,CarhartHarris2014} and therefore represents \emph{consistency} with the theory's commitments, not prediction-then-confirmation in the strict sense. The genuinely post-publication confirmations --- ConCrit \citep{AlgomShriki2026}, \citet{HengenShew2025}, \citet{Bhatt2024}, \citet{Li2025}, \citet{Pinto2017} --- constitute the stronger evidential claim. The distinction matters: consistency is evidence of non-refutation; post-publication convergence from independent groups is evidence of predictive validity.

\textbf{Anesthetic convergence on criticality disruption.} The theory predicts that all agents abolishing consciousness do so by pushing the substrate below the criticality threshold, regardless of receptor mechanism \citep{Gruber2015}. The Perturbational Complexity Index (PCI) threshold of 0.31 reliably discriminates conscious from unconscious states with high accuracy across propofol, midazolam, xenon, and ketamine \citep{Casali2013,Casarotto2016} --- work that predates the theory's publication but is consistent with its criticality commitment. Ketamine, which preserves criticality markers, preserves consciousness despite profound pharmacological effects. The ConCrit framework \citep{AlgomShriki2026} and the criticality meta-analysis by \citet{HengenShew2025} have independently reached the same conclusion from empirical consolidation of over 140 datasets. Tucker, Luu, and \citet{TuckerLuuFriston2025} provide independent neurobiological grounding, demonstrating that consciousness depends on excitatory-inhibitory balance maintaining cortical networks at the critical transition --- precisely the operating regime in which the four models interact. Their dual limbic architecture (predictive dorsal / corrective ventral) maps onto the implicit-to-explicit modeling axis. Recent work using adversarial AI trained on over 680,000 electrophysiological recordings identifies excessive cortical inhibitory coupling and basal ganglia indirect pathway disruption as the circuit-level mechanisms of unconsciousness in disorders of consciousness \citep{Toker2026}, consistent with the prediction that consciousness loss tracks a subcritical push.

\textbf{Waking criticality decline and sleep-dependent restoration.} The theory's sleep architecture prediction --- that waking degrades criticality and sleep restores it --- has been confirmed by \citet{Bhatt2024}, who demonstrated in continuous 10--14 day recordings that normal waking experience progressively disrupts criticality, and that sleep restores optimal computational regime. \citet{Meisel2013} showed fading criticality signatures during sustained human wakefulness. The theory predicted this from the principle that an analog substrate cannot sustain digital computation indefinitely without periodic recalibration.

\textbf{Sleep onset as bifurcation.} The prediction that sleep onset should be a radical transition (CA breakdown) rather than gradual dimming is supported by \citet{Li2025}, who demonstrated in over 1,000 participants that falling asleep follows a predictable bifurcation dynamic --- a tipping point preceded by critical slowing (increased variance and autocorrelation), with the transition detectable approximately 4.5 minutes before conventional sleep onset.

\textbf{Psychedelic content maps the processing hierarchy.} The prediction that increasing permeability exposes intermediate processing stages in hierarchical order is consistent with the independently developed REBUS model \citep{CarhartHarris2019}, \citet{Kluver1966} form constants, \citeauthor{Bressloff2002}\textquotesingle s (2002) V1 mathematical models, and dose-dependent visual cortex effects demonstrated with psilocybin \citep{Timmermann2023}. However, since predictive processing frameworks generate a nearly identical prediction through relaxation of top-down priors, this convergence validates the permeability principle without uniquely distinguishing the Four-Model Theory.

\textbf{Independent theoretical convergence on core architecture.} The theory's load-bearing features --- criticality, self-reference, holographic encoding, and autopoietic closure --- have received independent convergent support from two frameworks arriving from different methodological traditions. \citet{Bieberich2026} RIFT (Recurrent Integration Fractal Theory) framework, arising from fractal neuroscience, posits consciousness as a self-referential holographic endospace sustained at criticality through autopoietic feedback --- agreement on all four structural features with no terminological overlap and 25 years of independent development. Laukkonen, Friston, and \citet{Laukkonen2025} ``Beautiful Loop'' framework, arising from active inference, proposes that consciousness is constituted by a self-evidencing loop in which the system's generative model includes a model of itself generating the model --- converging on the self-referential closure mechanism that FMT identifies as the constitutive feature distinguishing conscious from non-conscious computation (Section 3.4.3). That three independent frameworks --- the Four-Model Theory from computational universality, RIFT from fractal dynamics, and Beautiful Loop from variational inference --- converge on self-referential closure at criticality as the constitutive mechanism of consciousness strengthens the case that these are constraints on any adequate theory rather than artifacts of a particular formalism.

\textbf{Split-brain holographic degradation.} The prediction that callosotomy produces bilateral degradation rather than clean hemispheric specialization is consistent with \citeauthor{Pinto2017}\textquotesingle s (2017) finding of ``unified consciousness, split perception'' --- each hemisphere retains a degraded but functionally complete conscious agent. The observable prediction is supported by recent evidence; the holographic storage mechanism that the theory proposes as explanation remains a distinctive theoretical contribution.


\subsection{Prediction 1: Psychedelics Alleviate Anosognosia}


\textbf{Statement.} Administration of psychedelic substances at sub-ego-dissolution doses to patients with anosognosia should produce a measurable increase in explicit awareness of their deficits. The effect should be dose-dependent and temporary, correlating with global increases in neural complexity.

\textbf{Mechanism.} \emph{Derivation: Variable implicit-explicit permeability (Principle 4) predicts that any manipulation changing global permeability should interact with any condition involving local permeability change. The prediction follows from the single-principle account: if permeability is one mechanism with local and global modes, a global increase should compensate for a local block.} Anosognosia is a local decrease in implicit-explicit permeability: the ISM registers the deficit, but the transfer to the ESM is blocked for that specific domain. Psychedelics globally increase permeability. The theory predicts that this global increase will compensate for the local block, allowing deficit information to reach the conscious simulation. This is the exact inverse of the psychedelic mechanism described in Section 6.1: psychedelics make the implicit accessible, anosognosia makes it inaccessible, and the two should interact predictably.

\textbf{Testability.} Clinical investigation of psychedelic effects on anosognosia patients --- beginning with case studies or observational reports, progressing to controlled trials. The prediction is specific: improvement in explicit awareness (measurable via established anosognosia scales), dose-dependent, and correlated with EEG complexity increases over the lesioned hemisphere. Psilocybin clinical trials already have extensive safety data at the relevant dose ranges.

\textbf{Falsification.} If psychedelics increase global neural complexity without improving explicit awareness in anosognosia, the variable-permeability model is incomplete. If the effect is not dose-dependent, the mechanism is different from what the theory proposes.

\textbf{Distinguishing power.} No consciousness theory predicts a psychedelic--anosognosia connection. This is a cross-domain surprise prediction: it connects psychopharmacology and clinical neurology through a single principle (variable permeability) in a way that no other framework generates. IIT, GNW, HOT, PP, and AST are all silent on why a psychedelic should help an anosognosia patient recognize their paralysis.


\subsection{Prediction 2: Ego Dissolution Content Is Controllable}


\textbf{Statement.} During psychedelic ego dissolution, the content of the altered identity experience --- what the subject reports ``being'' or ``becoming'' --- tracks the dominant sensory input. By controlling the sensory environment during ego dissolution, the identity content can be predicted and directed.

\textbf{Mechanism.} \emph{Derivation: The redirectable ESM (Principle 3) combined with variable permeability (Principle 4) predicts that when permeability increases beyond the point where normal self-referential input is disrupted, the ESM --- which requires input to function --- must latch onto whatever input remains dominant. Control the dominant input, control the identity content.} The ESM, deprived of normal self-referential input during ego dissolution, latches onto whatever input dominates the available stream. This is already suggested by salvia divinorum phenomenology, where users reliably report ``becoming'' objects in their immediate environment (furniture, walls, television characters). The theory predicts that this is not random but systematic: control the input, control the identity content. The prediction can be tested by varying the dominant sensory modality (visual, auditory, tactile) during ego dissolution and measuring correspondence between controlled input and reported identity content using blinded raters and established instruments such as the Ego Dissolution Inventory \citep{Nour2016}.

\textbf{Falsification.} If identity content during ego dissolution shows no systematic relationship to the sensory environment --- if subjects report identity transformations uncorrelated with the dominant input --- then the redirectable-ESM mechanism is wrong.

\textbf{Distinguishing power.} No competing theory generates this prediction. IIT, GNW, HOT, and AST have no mechanism for specifying what a subject will ``become'' during ego dissolution. Predictive processing (REBUS) predicts self-model relaxation but not the specific input-tracking pattern --- it explains \emph{that} the self dissolves, not \emph{what replaces it}. This is arguably the theory's most distinctive empirical prediction.


\subsection{Prediction 3: DID Alter Switches Show Self-Referential Network Specificity}


\textbf{Statement.} In dissociative identity disorder, controlled alter switching will produce neural reconfiguration patterns in which the representational dissimilarity (cosine distance between multi-voxel activation vectors) between alter states is significantly greater in self-referential processing regions than in sensorimotor regions (predicted cosine distance d $\geq$ 0.5 in self-referential regions vs. d < 0.2 in sensorimotor regions). These patterns will be alter-specific, reproducible across sessions, and distinguishable from the patterns produced by role-playing or method acting in matched controls.

\textbf{Mechanism.} \emph{Derivation: Simulation forking (Principle 5) predicts that multiple ESM configurations can run on a single substrate. Since forking operates on the ESM specifically --- not on the EWM or implicit models --- the neural differences between forked configurations should be concentrated where the ESM is computed, while leaving world-model and sensorimotor processing invariant.} Each alter is a distinct ESM configuration running on the same substrate with different parameters. The four-model architecture predicts that identity-level differences will cluster in networks subserving self-referential processing --- regions associated with self-narrative, autobiographical memory, and body ownership --- while leaving sensory and motor representations relatively invariant across alters. The prediction is formulated in terms of representational dissimilarity rather than anatomical localization to avoid the naive-modularism pitfall of equating the ESM with any specific brain network: the relevant signature is a \emph{gradient} of alter-specificity across the self-referential-to-sensorimotor axis, not activation of a named region. This goes beyond existing DID neuroimaging findings \citep{Reinders2003,Reinders2008,Schlumpf2014}, which demonstrate alter-specific activation differences but have not tested whether those differences follow the self-referential gradient the theory predicts.

\textbf{Falsification.} If representational dissimilarity between alter states is uniformly distributed across self-referential and sensorimotor regions --- if there is no gradient --- then the ESM-specificity claim is wrong. If whole-brain classification of alter identity does not significantly outperform sensorimotor-region-only classification, the prediction fails.

\textbf{Distinguishing power.} IIT predicts alter differences in posterior cortex integration. GNW predicts differences in prefrontal ignition patterns. Predictive processing predicts distributed differences in hierarchical prediction error. Only the Four-Model Theory predicts that representational dissimilarity should follow a self-referential gradient, because only FMT identifies each alter as a distinct ESM configuration whose differences should be concentrated where the self-model is computed.


\subsection{Prediction 4: Lucid Dream Onset Is a Criticality Threshold Crossing}


\textbf{Statement.} The transition from non-lucid to lucid dreaming during REM sleep corresponds to a threshold crossing in the ESM's criticality state, originating in self-referential cortical regions before spreading to other regions. The theory predicts that this transition has the structural signature of a phase transition --- critical slowing beforehand, discontinuous onset --- rather than a gradual ramp.

\textbf{Mechanism.} \emph{Derivation: Criticality (Principle 1) combined with the graduated consciousness hierarchy (Section 3.5) predicts that ESM activation depth varies with the substrate's position relative to the critical point. During REM, the substrate approaches but does not fully reach waking criticality; lucid onset corresponds to the ESM crossing from basic to simply extended consciousness --- a threshold, not a gradient.} During REM sleep, the substrate periodically re-approaches the criticality threshold as part of the NREM/REM restoration cycle. Lucid dreaming occurs when the substrate reaches sufficient criticality for the ESM to activate more fully --- a threshold crossing that produces the characteristic self-aware ``I am dreaming'' experience. The ESM activation should manifest as a criticality increase originating in self-model networks. This can be tested using the established lucid-dreamer eye-signaling paradigm \citep{LaBerge1985} with concurrent high-density EEG.

A methodological caveat is warranted. The specific criticality markers best suited to detect this transition --- whether Lempel-Ziv complexity, neuronal avalanche exponents, long-range temporal correlations, or measures not yet developed --- cannot be specified with confidence until the field has established reliable real-time tracking of criticality signatures during waking consciousness. Existing work on lucid dream onset \citep{Voss2014} reports gradual gamma-power increases rather than step-like transitions, but gamma power is a spectral measure, not a criticality measure; the two may dissociate. The prediction is structural (phase-transition dynamics in the ESM), not tied to any particular EEG frequency band. Its testability therefore depends on methodological advances in criticality measurement during sleep that are underway but not yet mature.

\textbf{Falsification.} If purpose-built criticality measures (branching ratio, DFA exponent, avalanche statistics) show no discontinuity at verified lucid onset --- or if the criticality change originates in sensory cortices rather than self-referential regions --- the prediction fails. The prediction is not falsified by gradual changes in spectral power, which track a different property.

\textbf{Distinguishing power.} IIT predicts increased $\Phi$ in the posterior hot zone --- a different spatial prediction. GNW predicts prefrontal ignition, which overlaps with the mPFC prediction but without the criticality-threshold framing. Predictive processing predicts increased precision on self-model predictions --- compatible but does not predict a phase-transition signature. The combination of phase-transition dynamics + ESM-network origination + criticality markers is unique to the Four-Model Theory.


\subsection{Theoretical Implication 1: Criticality + Four Models = Consciousness in Artificial Substrates}


\textbf{Statement}: A synthetic system implementing the four-model architecture at criticality will exhibit consciousness. The qualitative difference between interacting with such a system and interacting with a current LLM will be immediately and qualitatively distinguishable.

\textbf{Mechanism}: Substrate independence. Consciousness depends on function (four models at criticality), not on material (biological neurons).

\textbf{Testability}: Low at present (requires significant engineering development). However, intermediate tests are possible: systems with partial implementations (e.g., two models instead of four, or four models without criticality) should show partial consciousness indicators, detectable through behavioral signatures and neural/computational complexity measures. This implication is listed separately from Predictions 1--4 because it is not testable with current methods.

\textbf{Unique?}: Yes in the specific form. Other theories (PP, GNW) are compatible with artificial consciousness but do not provide a specific architectural blueprint.


\subsection{Theoretical Implication 2: Building Consciousness}


If the Four-Model Theory is correct, it should be possible to \emph{build} a conscious machine by implementing the specified architecture: four nested models (IWM, ISM, EWM, ESM) operating at criticality on a substrate of sufficient complexity. The theory implies that such a system would not merely \emph{imitate} consciousness (as an LLM might be said to do) but would \emph{be} conscious --- would have genuine phenomenal experience constituted by its virtual models.

This implication is not currently testable --- the engineering does not yet exist, and even if it did, the other-minds problem would make verification philosophically difficult. The other-minds problem applies universally: no behavioral test can conclusively demonstrate consciousness, whether in biological or artificial systems. This limitation is shared by all theories of consciousness, not specific to the Four-Model Theory. However, the implication sets a bold bar: if the theory is correct, the difference between interacting with a conscious artificial system and interacting with any current AI should be qualitatively obvious to human observers, just as the difference between a conversation with a conscious human and a conversation with a cleverly programmed chatbot is (according to the theory) a difference in \emph{kind}, not merely in degree.



%=============================================================================
\section{Open Questions}
%=============================================================================


Intellectual honesty requires identifying what the theory does not yet resolve. These are research frontiers, not theoretical weaknesses --- they are questions that arise \emph{from} the theory and that the theory's framework helps to sharpen.

\textbf{1. Are all four models virtual?} The theory as presented describes the implicit models (IWM, ISM) as ``real side'' and the explicit models (EWM, ESM) as ``virtual side.'' But it is not certain that this division is sharp. The implicit models might themselves have virtual properties --- they are, after all, \emph{models}, not raw physics. If the implicit models are also virtual in some sense, what constitutes the ``real side''? The raw physical substrate with no model-level description? This is an open question even for the theory's author, and its resolution may have consequences for the theory's treatment of the Hard Problem.

\textbf{2. Mathematical formalization.} The theory's criticality requirement is specified qualitatively (Wolfram's Class 4 regime), not quantitatively. A full mathematical treatment --- defining the four models formally, specifying the criticality threshold in terms of measurable quantities, deriving the predictions as formal consequences --- remains to be developed. The ConCrit framework's mathematical tools (power-law exponents, detrended fluctuation analysis, branching parameters) provide a starting point, as do the formal tools of dynamical systems theory, information geometry, and the category-theoretic approach to mathematical structure in conscious experience developed by \citet{Kleiner2024}.
\textbf{3. Physical implementation.} Which physical mechanism in the biological brain supports criticality? Candidates include cortical column dynamics, thalamocortical standing waves, glial modulation, and (more speculatively) quantum processes in microtubules \citep{PenroseHameroff1994}. The Four-Model Theory is agnostic here: it specifies the \emph{functional} requirements without mandating a specific physical mechanism. Resolving this question is an empirical task for neuroscience.

\textbf{4. Minimum configuration for consciousness.} Can the four models partially dissociate? Is it possible to have an EWM without an ESM (world-experience without self-experience), or an ESM without an EWM (self-experience without world-experience)? What is the minimum set of models required for consciousness, versus self-aware consciousness, versus full human-type consciousness? The theory's graduated levels (Section 3.5) suggest a hierarchy, but the exact minimum configuration may require simulation or empirical investigation to determine.

\textbf{5. Multi-level substrate architecture.} The biological brain can be analyzed as a hierarchy of nested systems: the physical substrate; within it, the proteomic (molecular) and topological (connectivity) systems, which mutually constrain each other; arising from these, the electrochemical system (neural signaling dynamics); and emerging from that, the virtual system --- the cortical automaton that hosts the four models. Each level both shapes and is shaped by its neighbors: the virtual system, over time, modifies the topological structure of the substrate it runs on. This raises a question for artificial consciousness: which levels of this hierarchy are essential, and which are specific to biological implementation? The theory's substrate-independence claim (Section 4.4) implies that only the virtual level is strictly required, but the bidirectional causal flow between levels suggests that decoupling the virtual system from its lower-level supports may not be straightforward.

\textbf{6. Decoding the virtual side.} The theory's real/virtual distinction has a methodological consequence that deserves explicit statement: the virtual side --- the conscious simulation --- is not directly accessible from external observation. Neuroimaging techniques (fMRI, EEG, MEG) capture substrate-level activity: blood-oxygen-level-dependent signals, electrical field potentials, magnetic flux. These are real-side measurements. They correlate with conscious states but do not constitute a readout of the simulation itself. To actually decode the content of conscious experience from neural data would require understanding the ``programming language'' of the brain --- the mapping from substrate dynamics to virtual content --- which in turn would require something approaching a full simulated connectome. The theory tells us \emph{what} the simulation is and \emph{where} it lives (in the explicit models, at the virtual level), but it does not provide the decoder ring. Developing this decoder --- bridging from substrate measurement to virtual-content readout --- constitutes a concrete research programme that follows directly from the theory's architecture and would, if achieved, provide the most direct possible test of the real/virtual distinction.

\textbf{7. The holography-criticality nexus.} The theory invokes both holographic storage (Section 5.2) and Class 4 criticality (Section 3.7) as essential features of the conscious substrate. The formal relationship between these two properties remains unexplored and presents a challenge suited to mathematical and computational investigation. Three distinct possibilities merit examination:

(a) \textbf{Holographic substrate $\rightarrow$ Class 4 dynamics.} Does a neural substrate that stores information holographically (in the patchwork sense defined in Section 5.2) necessarily exhibit Class 4 dynamics under appropriate driving conditions? If so, criticality would be a \emph{consequence} of the storage architecture rather than an independent requirement, simplifying the theory's axioms.

(b) \textbf{Class 4 automaton with holographic rule structure.} Can a cellular automaton be constructed whose transition rules are themselves defined holographically --- where each cell's update rule is a distributed function of the global rule set, degrading gracefully under partial rule deletion? Such a system would unify the holographic and criticality principles at the level of the automaton's definition rather than its behavior, and its characterization would constitute progress toward a mathematical formalization of the theory (see Open Question 2).

(c) \textbf{Class 4 dynamics $\rightarrow$ holographic emergent behavior.} Does a Class 4 automaton, regardless of its rule structure, necessarily produce holographic-like properties in its emergent patterns --- distributed information, graceful degradation, part-contains-whole structure? If this can be demonstrated, the holographic character of neural information storage would be derivable from the criticality requirement alone. A system satisfying both (a) and (b) --- holographic rules generating holographic output through Class 4 dynamics --- would constitute a computational fixed point: a process that encodes its own structure at every level. If such a system exists, it would suggest a deep connection between critical computation, holographic encoding, and self-referential structure --- potentially relevant to cosmological questions about the computational nature of physical law (see Gruber, 2026c for a detailed treatment).

Resolving which of these relationships holds (and whether multiple hold simultaneously) is a mathematical challenge that lies at the intersection of cellular automata theory, information theory, and the theory of distributed computation. These questions are speculative and lie at the boundary of the theory's current scope, but they are well-defined mathematical problems. The author invites researchers in these fields to investigate these conjectures, which may bear on questions in computational complexity well beyond consciousness science.



%=============================================================================
\section{Discussion}
%=============================================================================



\subsection{Implications for Artificial Consciousness}


The Four-Model Theory provides an engineering specification for artificial consciousness: implement the four-model architecture on a substrate operating at criticality. This is a concrete deliverable, not an abstract philosophical claim.

Current AI systems fail this specification in at least two ways. Large language models (LLMs) operate via feedforward inference (transformer attention is computed in a single pass without recurrent dynamics), which corresponds to Wolfram's Class 1/2 --- far below the criticality threshold. They also lack the four-model architecture: there is no ISM (no substrate-level self-knowledge that is \emph{distinct from} the model's outputs), no ESM (no ongoing self-simulation that constitutes a subjective perspective), and no real/virtual split (no two-level ontology in which experience resides at the virtual level).

This does not mean that LLMs are necessarily non-conscious --- the theory cannot prove a negative --- but it predicts that they lack the architecture required for consciousness as the theory defines it. The growing discourse around AI consciousness \citep{Butlin2023,Butlin2025,Schwitzgebel2025,Birch2025} and AI welfare \citep{Long2024,Anthropic2025} makes this distinction practically important. A theory that provides clear criteria for artificial consciousness --- rather than vague analogies to human cognition --- has immediate ethical and engineering value. The intelligence implications of this architectural specification are explored in \citet{Gruber2026a}, which argues that current AI systems fail to exhibit self-directed intellectual development precisely because they lack the motivational component that drives the recursive intelligence loop.


\subsection{Implications for Consciousness Science}


The Four-Model Theory suggests a shift in experimental priorities. Rather than adjudicating between IIT and GNW (the current focus of adversarial collaborations), the field should:

\begin{enumerate}
\item \textbf{Investigate the psychedelic-anosognosia connection} (Prediction 1) --- a cross-domain surprise prediction that, if confirmed, would constitute strong evidence for the variable-permeability mechanism.
\item \textbf{Design controlled ego-dissolution experiments} (Prediction 2) --- the most distinctive prediction, uniquely generated by the redirectable ESM mechanism.
\item \textbf{Test ESM-network localization in DID} (Prediction 3) --- combining established DID neuroimaging paradigms with network-specific analysis to test whether alter differences are DMN-concentrated.
\item \textbf{Measure criticality at lucid dream onset} (Prediction 4) --- achievable with established lucid-dreamer paradigms and high-density EEG.
\end{enumerate}

These experiments do not require committing to the Four-Model Theory in its entirety. They test specific mechanisms (criticality, redirectable ESM, variable permeability) that could be incorporated into other frameworks if confirmed.


\subsection{Limitations}


\textbf{No institutional laboratory.} The predictions presented here were derived theoretically and have not been tested in the author's own laboratory. While this does not affect their validity as predictions, it does mean that empirical testing depends on the willingness of established laboratories to take them up.

\textbf{The causal status of experience is nuanced and will draw fire from both sides.} The theory holds that qualia lack independent causal power over the substrate, yet the simulation of which they are constitutive is functionally essential --- the substrate's own mechanism for consequence-evaluation (Section 4.2). This is distinct from classical epiphenomenalism, in which consciousness is a causally inert by-product. The position will face resistance from two directions: philosophers who insist that consciousness must have independent causal efficacy will reject the denial of direct causal power, while strict epiphenomenalists will object that the feedback loop between explicit and implicit models smuggles causation back in under another name. The defense offered in Section 4.2 is, I believe, internally consistent, but the reader should be aware that the theory's causal architecture occupies contested ground between these camps.

\textbf{Qualitative rather than quantitative.} The theory's predictions are currently stated in qualitative terms (``criticality increases,'' ``permeability changes,'' ``ESM redirects''). Quantitative formalization would strengthen the predictions and enable more precise experimental testing. This is acknowledged as a priority for future work.

\textbf{The other-minds problem.} The ultimate prediction --- that a system built to the theory's specification would be conscious --- faces the standard other-minds problem: how would we verify consciousness from the outside? The theory predicts that the difference would be qualitatively obvious to human observers, but ``qualitatively obvious'' is not a measurement. Developing consciousness indicators that can be applied to artificial systems is a challenge for the entire field, not specific to this theory.

\textbf{Inherent limits of inside-modeling.} The theory proposes a self-model as the basis of consciousness --- but the modeler and the modeled are the same system. This raises a structural concern analogous to G{\"o}del's incompleteness: a formal system cannot prove all truths about itself. Similarly, the brain attempting to model itself faces an irreducible epistemological gap --- the instrument is the object of study. The theory acknowledges this limitation explicitly through the Meta-Problem (Section 3.8): the ESM cannot observe the ISM's mechanisms, which is why consciousness seems mysterious \emph{from the inside}. This is a feature of the architecture, not a flaw in the theory, but it does imply that any theory of consciousness generated by a conscious system will carry a residual blind spot.

\textbf{Language linearizes non-linear phenomena.} Any theoretical framework expressed in natural language necessarily serializes phenomena that may be inherently parallel and non-linear. The four models operate simultaneously in high-dimensional state spaces; describing them sequentially in prose introduces a representational loss that no amount of careful writing can fully eliminate. This limitation applies to all theories of consciousness, not only the present one, but it should be kept in mind when evaluating the theory's verbal descriptions against the multidimensional dynamics they attempt to capture.

\textbf{Criticality-rhythm relationship not formalized.} Section 5.1 argues that criticality provides punctuated stability, with biological rhythms (sleep-wake cycling, ultradian rhythms, neurotransmitter depletion and replenishment) governing how long the substrate can maintain the critical regime. This relationship between dynamical criticality and biochemical rhythm is proposed qualitatively; a formal model linking neurotransmitter kinetics to criticality maintenance and breakdown remains to be developed.

\textbf{Every model has modeling error.} The Four-Model Theory is itself a model --- and every model, by definition, is a simplification that carries inherent modeling error. The theory does not claim to be a final or complete account of consciousness; it claims to be a useful one, generating testable predictions and unifying phenomena that other frameworks address in isolation. To the extent that it is wrong, the predictions in Section 8 are designed to reveal where.



%=============================================================================
\section{Conclusion}
%=============================================================================


The Four-Model Theory of Consciousness proposes that consciousness is an ongoing self-simulation across four nested models --- Implicit World Model, Implicit Self Model, Explicit World Model, and Explicit Self Model --- operating on a substrate at the edge of chaos. Qualia are virtual: they are the phenomenal properties of the simulation, not of the substrate. The theory addresses the Hard Problem by revealing a category error in its formulation, simultaneously closing the Explanatory Gap and accounting for the Meta-Problem.

The theory addresses all eight requirements for a complete theory of consciousness: the Hard Problem (addressed via virtual qualia), the Explanatory Gap (addressed alongside), the Boundary Problem (defined by the scope of virtual models), the Structure of Experience (generated by the simulation's complexity), Unity and Binding (emergent from critical dynamics), Combination and Emergence (weak emergence, no combination problem), the Causal Role (the architecture is causally efficacious; qualia lack independent causal power but are constitutive of the simulation the substrate deploys for evaluation), and the Meta-Problem (structural inaccessibility of the ISM to the ESM).

The theory generates four novel testable predictions, each with clear falsification criteria: that psychedelics should alleviate anosognosia via the variable-permeability mechanism (Prediction 1), that ego dissolution content tracks dominant sensory input (Prediction 2), that DID alter switches are concentrated in ESM-related networks (Prediction 3), and that lucid dream onset is a criticality threshold crossing (Prediction 4). All four are distinctive to the Four-Model Theory. Additionally, several claims originally derived from the theory's axioms --- including the anesthetic-criticality convergence and sleep-dependent criticality restoration --- have been independently confirmed by subsequent empirical work (Section 8.1).

The theory's criticality requirement was derived from Wolfram's computational framework in 2015, independently of the empirical criticality program already underway \citep{Beggs2003,CarhartHarris2014}. The later large-scale consolidation --- \citet{HengenShew2025} meta-analysis of 140 datasets and \citet{AlgomShriki2026} ConCrit framework --- confirmed the criticality-consciousness link. This convergence --- a theoretical prediction derived independently from computational first principles, later confirmed by large-scale empirical synthesis --- provides notable support.

Open questions remain: the status of the implicit models (real or also virtual?), the need for mathematical formalization, the specific physical mechanism supporting criticality, and the minimum configuration for consciousness. These are research frontiers that the theory's framework helps to sharpen.

The ambition of consciousness science is not merely to correlate neural activity with subjective reports but to understand \emph{why} there is experience at all. The Four-Model Theory offers an answer: there is experience because there is a simulation, and within the simulation, experience is not an addition to the process but is constitutive of it. The way to test this answer is not through philosophical argument alone but through the predictions it generates --- and ultimately, through the engineering challenge of building a system to the specification and observing whether the result is, as the theory predicts, qualitatively unlike anything that exists today. The theory's implications extend beyond consciousness science: the cognitive-learning capacity it identifies as a consequence of the four-model architecture is the foundation for a recursive model of intelligence \citep{Gruber2026a} in which knowledge, performance, and motivation form a self-reinforcing loop.



\section*{Acknowledgments}


This paper was written with the assistance of Claude \citep{Anthropic2025}, which served as editor, cross-checker, and writing tool throughout the drafting process. The theory, arguments, predictions, and all intellectual content are entirely the author's, originally published in \citet{Gruber2015}; Claude assisted with literature cross-referencing, prose drafting from the author's directions, and internal consistency checking.

\section*{Data Availability}


No new data were generated or analysed in support of this research.

\section*{Funding}


This research received no specific grant from any funding agency in the public, commercial, or not-for-profit sectors.

\section*{Declaration of generative AI and AI-assisted technologies in the manuscript preparation process}


During the preparation of this work the author used Claude (Anthropic) in order to conduct structured adversarial challenges against the theory, assist with literature search and citation verification, and refine editorial presentation. The underlying theory was developed in 2015 without AI assistance, predating current generative AI capabilities. After using this tool, the author reviewed and edited the content as needed and takes full responsibility for the content of the published article.

\bibliographystyle{plainnat}

\bibliography{references}

\end{document}
